% main.tex -- H100 four-engine study (CPU/GPU Polars, CPU/GPU Spark).
%
% DRAFT. Methods written before any timed data existed; results prose drafted
% from the final analysis outputs (2026-10-30). Every result is a \stat{} macro
% filled from the analysis output (tools/stats_to_tex.py), and every claim's
% strength follows its interval. The methods must match
% docs/H100_EXPERIMENT_DESIGN.md (locked 2026-10-29, amendments A1--A4); any
% further amendment must be reflected here.
%
% Host-conditional text (amendments A2, A3): the A100, GH200 and B200
% campaigns each either ran or did not. Their text sits between
% "% BEGIN-RAN:<host>" and "% END-RAN:<host>" lines (or NOTRUN for the
% not-run case; <host> "any" = at least one of them ran) and inside the
% matching \if<flag>ran conditional from macros.tex. The final build refuses
% to proceed until every host's outcome is stated (paper/README.md).
\documentclass[11pt]{article}

\usepackage[T1]{fontenc}
\usepackage[utf8]{inputenc}
\usepackage{lmodern}
\usepackage[margin=1in]{geometry}
\usepackage{amsmath,amssymb}
\usepackage{booktabs}
\usepackage{multirow}
\usepackage{graphicx}
\usepackage[table]{xcolor}
\usepackage[numbers,sort&compress]{natbib}
\usepackage{caption}
\usepackage{subcaption}
\usepackage{enumitem}
\usepackage{xspace}
\usepackage[hidelinks]{hyperref}
\usepackage{url}
\usepackage{authblk}
\usepackage{seqsplit}  % line-breakable commit hashes and SHA-256 digests

% macros.tex -- statistics lookup, figure inclusion and notice macros.
%
% Numbers in the paper are never typed in. Every result is written as \stat{key},
% where key is a dotted stats.json key (e.g. c1.baseline.sf300.estimate).
% paper/tools/stats_to_tex.py turns <analysis-out>/stats.json into
% generated/stats.tex, which contains one \defstat{key}{value} per key and
% \statsfinaltrue. In final mode a missing key is a hard LaTeX error; in draft
% mode (no generated/stats.tex) the key is printed in a box so the gaps are visible.

\makeatletter
\newif\ifstatsfinal
\statsfinalfalse
\newcommand{\defstat}[2]{\expandafter\def\csname stat@#1\endcsname{#2}}
\newcommand{\statplaceholder}[1]{%
  \fcolorbox{red}{yellow!20}{\scriptsize\ttfamily\detokenize{#1}}}
\newcommand{\stat}[1]{%
  \ifcsname stat@#1\endcsname
    \csname stat@#1\endcsname
  \else
    \ifstatsfinal
      \PackageError{stats}{Missing statistic '#1' in generated/stats.tex}%
        {Re-run tools/stats_to_tex.py; the analysis output lacks this key.}%
    \else
      \statplaceholder{#1}%
    \fi
  \fi}
% "estimate [ci_low, ci_high]" for a ratio prefix such as c1.baseline.sf300
\newcommand{\ratioci}[1]{%
  \stat{#1.estimate}\,{\footnotesize[\stat{#1.ci_low},\,\stat{#1.ci_high}]}}
% "estimate x (95% CI ci_low--ci_high)" for running text
\newcommand{\ratiotext}[1]{%
  \stat{#1.estimate}$\times$ (95\%~CI \stat{#1.ci_low}--\stat{#1.ci_high})}
\makeatother

% GPU-generation ladder (amendments A2, A3). One flag per additional host;
% the H100 is primary and always present. Draft builds show every host's text
% with placeholders; stats_to_tex.py --not-run HOST=REASON writes
% \<flag>ranfalse, and \ladderranfalse when no additional host ran. Wrap
% host text in \if<flag>ran ... \fi inside "% BEGIN-RAN:<host>" /
% "% END-RAN:<host>" comment lines (NOTRUN for the \else case) so the key
% scanner knows which keys each case needs.
\newif\ifaonehundredran  \aonehundredrantrue   % a100
\newif\ifghtwohundredran \ghtwohundredrantrue  % gh200
\newif\ifbtwohundredran  \btwohundredrantrue   % b200
\newif\ifladderran       \ladderrantrue        % any of the three ran
\expandafter\def\csname hostflag@a100\endcsname{aonehundred}
\expandafter\def\csname hostflag@gh200\endcsname{ghtwohundred}
\expandafter\def\csname hostflag@b200\endcsname{btwohundred}

% Disclosure watermarks, set by generated/stats.tex: hand-filled _display values
% (stats_to_tex.py --allow-manual-overrides, or a draft) and synthetic stats.
% Printed diagonally on every page and in the title block.
\newif\ifmanualvalues
\newif\ifsyntheticstats

\IfFileExists{generated/stats.tex}{\input{generated/stats.tex}}{}

\newcommand{\watermarktext}{%
  \ifmanualvalues MANUAL VALUES\fi
  \ifmanualvalues\ifsyntheticstats\ \textperiodcentered\ \fi\fi
  \ifsyntheticstats SYNTHETIC -- NOT RESULTS\fi}
\newcommand{\watermarkpage}{%
  \put(0.5\paperwidth,-0.5\paperheight){\makebox(0,0){\rotatebox{45}{%
    \scalebox{4}{\sffamily\bfseries\textcolor{red!30}{\watermarktext}}}}}}
\AddToHook{shipout/background}{%
  \ifmanualvalues\watermarkpage\else\ifsyntheticstats\watermarkpage\fi\fi}
\newcommand{\watermarknotice}{%
  \ifmanualvalues\\\textcolor{red}{\textbf{\watermarktext}}%
  \else\ifsyntheticstats\\\textcolor{red}{\textbf{\watermarktext}}\fi\fi}

% Host-conditional values for tables: the value if the host's campaign ran,
% \hostnotrun otherwise. Keys: <host>.<key>, <host>.<prefix>.{estimate,...}
% and xhost.<host>.<rest>.{estimate,...}. Not for use inside \if...\fi text
% that may be skipped (they are not conditionals themselves).
\makeatletter
\newcommand{\hostnotrun}{---}
\newcommand{\ifhost@ran}[1]{\csname if\csname hostflag@#1\endcsname ran\endcsname}
\newcommand{\hoststat}[2]{\ifhost@ran{#1}\stat{#1.#2}\else\hostnotrun\fi}
% \hoststator{host}{key}{fallback}: the value if the host ran, else the fallback text.
\newcommand{\hoststator}[3]{\ifhost@ran{#1}\stat{#1.#2}\else#3\fi}
\newcommand{\hostratioci}[2]{\ifhost@ran{#1}\ratioci{#1.#2}\else\hostnotrun\fi}
\newcommand{\xhostratioci}[2]{\ifhost@ran{#1}\ratioci{xhost.#1.#2}\else\hostnotrun\fi}
\makeatother

% Figures are copied into figures/ by tools/sync_figures.py. Until then a
% labelled placeholder box keeps the layout and cross-references working.
\newcommand{\studyfig}[2]{%
  \IfFileExists{figures/#1.pdf}{\includegraphics[width=#2]{figures/#1.pdf}}{%
  \IfFileExists{figures/#1.png}{\includegraphics[width=#2]{figures/#1.png}}{%
  \fbox{\parbox[c][0.28\textheight][c]{0.9\linewidth}{\centering\ttfamily
    figure pending: \detokenize{#1}}}}}}

% Author-supplied detail still missing (affiliation, email, venue, ...).
\newcommand{\todoplaceholder}[1]{\textcolor{red}{[TODO: #1]}}

% Results-not-yet-available marker for draft prose (also flagged with
% "% TODO(results)" comments in the source).
\newcommand{\pending}[1]{\textcolor{gray}{[pending results: #1]}}

% TPC fair-use disclaimer. The template wording is TPC's (Fair Use of TPC
% Benchmarks quick reference, 2019-06-05; current Policies v6.20 sec. 8.1.5).
% The name of the derived work is an open author decision; see paper/README.md.
\newcommand{\workloadname}{H100 study workload}
\newcommand{\tpcdisclaimer}{The \workloadname{} is derived from the TPC-H
  Benchmark and as such is not comparable to published TPC-H Benchmark results,
  as the \workloadname{} results do not comply with the TPC-H Benchmark
  Specification.}
\newcommand{\chartnote}{\emph{Derived from TPC-H; not comparable to published
  TPC-H results (see Section~\ref{sec:scope}).}}

% Engine and contrast names
\newcommand{\eng}[1]{\texttt{#1}}
% Prose names for running text and tables; \id... for the code identifiers
% (arm ids in run configuration and the seeded order assignment).
\newcommand{\polarscpu}{Polars CPU}
\newcommand{\polarsgpu}{Polars GPU}
\newcommand{\sparkcpu}{Spark CPU}
\newcommand{\sparkgpu}{Spark GPU}
\newcommand{\idpolarscpu}{\eng{polars}}
\newcommand{\idpolarsgpu}{\eng{polars\_gpu}}
\newcommand{\idsparkcpu}{\eng{spark\_cpu}}
\newcommand{\idsparkgpu}{\eng{spark\_gpu}}
\newcommand{\q}[1]{\texttt{#1}}

\setlength{\emergencystretch}{3em}

% Title. The suffix appears only if at least one additional ladder host ran.
% TODO(author): choose a title once the ladder outcome is known. Options:
%  (a) [current] CPU Dataframes versus GPU Acceleration on a Single Host: A
%      Pre-registered, Randomized Comparison of Polars and Spark on a
%      TPC-H-Derived Workload, Replicated across GPU Generations
%  (b) [all four hosts ran] Does a Newer GPU Help? Polars and Spark, on the CPU
%      and the GPU, from A100 to B200: A Pre-registered Comparison
%  (c) [result-led; only if the CIs support it] When Does a GPU Beat a Fast
%      CPU Dataframe Engine? A Pre-registered Four-Engine Study on Four GPU
%      Generations
%  (d) [no additional host ran] keep (a) without the suffix.
\title{CPU Dataframes versus GPU Acceleration on a Single Host:\\
  A Pre-registered, Randomized Comparison of Polars and Spark\\
  on a TPC-H-Derived Workload\ifladderran, Replicated across GPU Generations\fi}
% TODO(author): target venue is pending (no ORCID is given).
\author{Daniel Ejiofor}
\affil{Independent Researcher}
\affil{\href{mailto:daniel@siriusmed.io}{\nolinkurl{daniel@siriusmed.io}}}
\date{Draft of \today\watermarknotice}

\begin{document}
\maketitle

\begin{abstract}
% Result sentences drafted from the final analysis outputs (2026-10-30);
% every number is a \stat{} macro and every claim follows its interval.
We compare Polars and Spark, each on the CPU and on the GPU (both GPU paths
use NVIDIA cuDF), on one NVIDIA H100 host. The workload
is TPC-H-derived (six adapted TPC-H queries, SF1--SF300; six complex queries,
SF30--SF300), not a TPC-H benchmark result. Its pre-registered design uses
four rounds in a Williams Latin square, timing-independent integrity gates and
95\% intervals across rounds. At SF300, GPU speedups over the same
framework's CPU engine were \ratiotext{c1.baseline.sf300} for Polars and
\ratiotext{c2.baseline.sf300} for Spark on the baseline suite, and
\ratiotext{c1.complex.sf300} and \ratiotext{c2.complex.sf300} on the complex
suite. The best GPU engine was faster than Polars CPU from
\stat{crossover.baseline.c4.scale} (baseline) and
\stat{crossover.complex.c4.scale} (complex), and slower at
\stat{c4.baseline.summary.below}. Polars GPU rejected
\stat{coverage.polars_gpu.complex.unsupported} of the six complex queries as
unsupported (CPU fallback disabled by design). However, the GPU engine
requires Polars 1.44.x, and Polars 2.0.0 on the CPU was
\stat{bridge_summary.both}. In an exploratory, post-hoc comparison with
Polars 2.0.0 (E1), the Polars GPU advantage shrank but stayed above 1
\stat{exploratory.e1.baseline.summary.above_phrase} (baseline).
% BEGIN-RAN:any
\ifladderran
Campaigns on \stat{hosts.other_ran_list} showed the same pattern, with GPU
time ratios to the H100 of
% BEGIN-RAN:a100
\ifaonehundredran
\stat{ladder.xhost_gpu.a100.min} to \stat{ladder.xhost_gpu.a100.max} on the
A100 (slower) and
\fi
% END-RAN:a100
\stat{ladder.xhost_gpu.later.min} to \stat{ladder.xhost_gpu.later.max} on the
\stat{ladder.xhost_gpu.later.hosts}.
\fi
% END-RAN:any
Across \stat{ladder.hosts_count} hosts, \stat{ladder.outputs_failed_total} of
\stat{ladder.outputs_compared_total} output comparisons disagreed with Polars CPU.
\end{abstract}

\noindent\textbf{Scope notice.} \tpcdisclaimer{} The queries are a selected
and adapted subset with additional custom queries, Float64 normalization of
decimal columns, and our own timing and summary method; see
Section~\ref{sec:scope}.

% ======================================================================
\section{Introduction}
\label{sec:intro}

Practitioners who outgrow a laptop-scale dataframe face two common routes. One
is a faster single-node engine: Polars~\citep{polars1442} executes lazy,
optimized query plans with a multi-threaded and streaming engine written in
Rust. The other is GPU acceleration, either inside the same dataframe API (the
Polars GPU engine built on NVIDIA cuDF~\citep{polars_gpu_guide,cudf_polars_docs})
or underneath an existing Spark SQL workload (NVIDIA's cuDF plugin for Apache
Spark~\citep{nvidia_cudf_spark_docs}). Vendor and project material reports
large speedups for both routes~\citep{semaan2024polarsgpu,nvidia_spark_benchmarks},
usually against a CPU baseline from the same framework.

The decision a practitioner actually faces is broader: \emph{for this kind of
query at this data size, is a GPU engine faster than the best CPU engine I could
use on the same machine?} Answering that requires holding the host, the input
files, the query semantics and the timing boundary fixed across all engines,
checking that every engine returns the same answer, and checking that a ``GPU''
timing really ran on the GPU. It also requires enough repeated, order-balanced
measurement to separate engine effects from run-to-run noise and order
effects~\citep{mytkowicz2009wrong,kalibera2013rigorous,hoefler2015scientific,raasveldt2018fair}.

An earlier study by the author compared three of these engines (Polars CPU,
Spark CPU, and Spark with the GPU plugin) on an A100 40\,GB host. It found that
the plugin accelerated Spark on the larger scales but that Polars CPU retained
the lower aggregate score at every observed scale, with GPU Spark winning only
selected queries. That study ran one session per engine and scale, so it
had no replication across sessions and no interval estimates, and it did not
include the Polars GPU engine. This paper extends it.

\paragraph{Contributions.}
\begin{enumerate}[leftmargin=*,itemsep=2pt]
  \item A four-engine comparison (CPU and GPU variants of Polars and Spark) on
    one H100 host with identical, verified inputs, from SF1 to SF300
    (1.8~billion \texttt{lineitem} rows).
  \item A pre-registered design~(Section~\ref{sec:design}): a Williams Latin
    square over four rounds, fixed contrasts, and a statistical plan that
    reports geometric-mean ratios with across-round confidence intervals rather
    than single-run point values.
  \item Integrity gates that are independent of timings: cross-engine output
    agreement, a strict no-CPU-fallback guard for the Polars GPU engine, and
    final-plan checks for the Spark plugin.
  \item A GPU coverage analysis that reports which queries cannot run on each
    GPU engine, rather than silently timing a CPU fallback.
  \item A side study measuring how much the Polars version (1.44.2 versus
    2.0.0) changes CPU timings, bridging to the earlier A100 study.
% BEGIN-RAN:any
\ifladderran
  \item A GPU-generation ladder: the complete campaign repeated, with the same
    design, software versions, inputs and engine thread count, on
    \stat{hosts.other_ran_list}, each analyzed separately and compared
    descriptively with the H100 (Section~\ref{sec:ladder}).
% BEGIN-RAN:a100
\ifaonehundredran
    On the A100 it supplies the four-engine result that the earlier A100
    study lacked.
\fi
% END-RAN:a100
\fi
% END-RAN:any
  \item A complete artifact: code, saved outputs, plans, telemetry and
    manifests (Section~\ref{sec:artifact}).
\end{enumerate}

% ======================================================================
\section{Background and Related Work}
\label{sec:background}

\paragraph{Dataframe engines.} The pandas library~\citep{mckinney2010pandas}
made the dataframe the default interface for tabular analysis in Python;
\citet{petersohn2020dataframes} discuss why its eager, single-threaded execution
model is hard to scale. Polars~\citep{polars1442,polars200} provides a lazy API
whose query plans are optimized (predicate and projection pushdown, among other
rewrites) and executed by a multi-threaded engine; its streaming engine
processes data in batches. DuckDB~\citep{raasveldt2019duckdb} is an embedded
analytical SQL engine; here it is used only to generate data.

\paragraph{GPU dataframes.} cuDF is NVIDIA's GPU dataframe library, built on the
C++ library libcudf. The Polars GPU engine (package \texttt{cudf-polars}) takes
an optimized Polars plan and executes it with cuDF; when a plan contains an
unsupported operation it falls back to the CPU engine unless the user sets
\texttt{raise\_on\_fail=True}~\citep{polars_gpu_guide,cudf_polars_docs}. The
Polars user guide states that the engine is in open beta, that speedups are
most likely for workloads dominated by grouped aggregations and joins, and that
I/O-bound queries perform about the same as on the CPU~\citep{polars_gpu_guide}.
Each \texttt{cudf-polars} release supports a bounded range of Polars versions;
this constraint shaped our version choice (Section~\ref{sec:arms}).

\paragraph{Spark and its GPU plugin.} Apache Spark generalizes MapReduce with
resilient distributed datasets~\citep{zaharia2012rdd,zaharia2016spark}; Spark
SQL adds a relational optimizer (Catalyst) and DataFrame
API~\citep{armbrust2015sparksql}. Adaptive query execution re-optimizes plans
at runtime using shuffle statistics~\citep{spark359docs}. NVIDIA's cuDF plugin
for Apache Spark~\citep{nvidia_cudf_spark_docs,nvidia_cudf_spark_repo}
replaces supported physical operators with GPU implementations through libcudf
and leaves the rest on the CPU, so a GPU-enabled Spark plan can be hybrid. We
therefore inspect final plans rather than assume GPU execution.

\paragraph{TPC-H and derived use.} TPC-H is the Transaction Processing
Performance Council's decision-support benchmark~\citep{tpch301,poess2000tpc}.
\citet{boncz2014tpch} analyze which ``choke points'' its 22 queries exercise.
Many engine evaluations reuse its schema, generator and queries without
following the full specification, under names that mark the derivation: the
Polars project's PDS-H~\citep{vink2025pdsh} and NVIDIA's NDS-H for the Spark
plugin~\citep{nvidia_spark_benchmarks} both state that their results are
derived from TPC-H and are not comparable to published TPC-H results. We
generate data with DuckDB's \texttt{tpch} extension, which documents small
differences from the official specification and fixed substitution
parameters~\citep{duckdb_tpch_ext}. TPC's fair-use material requires such
studies to carry a prescribed disclaimer, disclose deviations and the test
environment, and avoid comparison with official
results~\citep{tpcfairuse2019,tpcpolicies620}.

\paragraph{GPU and CPU analytics.} \citet{bress2014gpudb} and
\citet{rosenfeld2022heterogeneous} survey GPU-accelerated and heterogeneous
CPU/GPU query processing. \citet{shanbhag2020gpucpu} argue that because
analytical operators are usually memory-bandwidth bound, GPU speedups over
well-implemented CPU code should be expected near the GPU-to-CPU bandwidth
ratio rather than orders of magnitude, and analyze when and why measured
speedups depart from it. \citet{lutz2020pump} show that the CPU--GPU transfer
link bottlenecks GPU query processing on data larger than device memory, and
that faster interconnects relieve it. Public GPU dataframe comparisons are
mostly vendor or project benchmarks: the Polars GPU announcement reports
speedups over Polars CPU on PDS-H using an H100~\citep{semaan2024polarsgpu},
and the long-running database-like ops benchmark~\citep{dbbenchmark} compares
CPU tools on grouping and joins. To our knowledge, independent, replicated
comparisons of a GPU dataframe engine and a GPU Spark plugin against a strong
single-node CPU engine on identical inputs remain scarce.
% TODO(author): extend with any further related work the author knows of;
% candidates that could not be verified are listed in CITATIONS_TODO.md.

\paragraph{Benchmarking methodology.} Hidden factors such as environment
layout can bias measurements~\citep{mytkowicz2009wrong}; repeated, randomized
execution with interval estimates is recommended~\citep{kalibera2013rigorous,
hoefler2015scientific}, and database comparisons have many fairness
pitfalls~\citep{raasveldt2018fair}. Ratios of times should be summarized with
the geometric mean~\citep{fleming1986geomean}. Williams designs balance
first-order carryover in crossover experiments~\citep{williams1949latin}.

% ======================================================================
\section{Experimental Design}
\label{sec:design}

\subsection{Pre-registration and amendments}
\label{sec:prereg}

The research questions, arms, workload, order design, integrity gates,
pre-declared GPU gaps and statistical plan below were written down and locked
in the study repository (\texttt{docs/H100\_EXPERIMENT\_DESIGN.md}) on
2026-10-29, before any timed data existed on any host. The lock rule is that
the locked sections are never edited; any change is appended as a dated
amendment with its reason. Four amendments were made, A1--A3 before any
timed data existed on any host and A4 before any B200 data existed:
\begin{description}[leftmargin=*,itemsep=2pt]
  \item[A1 (failure classification, 2026-10-29).] As first written, a
    deterministic GPU resource exhaustion (for example, running out of device
    memory at SF300) would have failed a round and stopped the campaign, even
    though RQ4 treats such failures as results. A1 separates \emph{coverage
    outcomes}, which are reported, from \emph{integrity failures}, which fail
    a round (Sections~\ref{sec:gates} and~\ref{sec:coverage-design}).
  \item[A2 (B200 host, 2026-10-29).] A second host with an NVIDIA B200 runs
    its own complete campaign under the unchanged design. The H100 campaign
    remains the primary analysis (Sections~\ref{sec:hosts}
    and~\ref{sec:stats}).
  \item[A3 (GPU-generation ladder, 2026-10-30).] Two further hosts, an A100
    and a GH200, each run their own complete campaign exactly as A2 specifies,
    so the H100 result sits on a ladder from Ampere through Hopper and
    Grace~Hopper to Blackwell, and the four-engine result missing from the
    earlier A100 study is supplied. All hosts use 26 engine threads
    (Section~\ref{sec:ladder-design}).
  \item[A4 (known Blackwell compiler defect, 2026-10-30).] Checking the B200
    precondition found an open defect in the CUDA 12.9 \texttt{ptxas}
    compiler on Blackwell that can make
    \texttt{thrust::exclusive\_scan\_by\_key} return wrong results silently
    beyond its first tile~\citep{cccl11167}. The \texttt{libcudf} build used by
    \polarsgpu{} contains cuDF's workaround~\citep{cudf23973}; the plugin JAR
    used by \sparkgpu{} predates it. A4 changes no gate and no analysis: on the
    B200, an output disagreement remains an integrity failure, and its diagnosis
    records whether the evidence fits the defect (Section~\ref{sec:threats}).
\end{description}
No decision about the design, stopping or analysis was based on observed
speeds.

\paragraph{Timestamp.} The locked design together with amendments A1--A3 was
committed as \texttt{ee2e6a7} (full hash
\texttt{\seqsplit{ee2e6a76a0f0b8ddd28bc4a9293dde5176c34ffe}}) on branch
\texttt{H100\_Test} of
\url{https://github.com/ejiod2025/gpu-spark-polars-study} and pushed on
2026-10-30 (commit time 2026-10-30 03:55\,UTC). The design file at that commit
has SHA-256
\texttt{\seqsplit{6b7544ee23b4ad40d66c8ba0866772ce0fec0ca6a6db92b6435f064979dcddff}};
readers can check it with \texttt{git show
ee2e6a7:docs/H100\_EXPERIMENT\_DESIGN.md | sha256sum}. The lock and
amendments A1--A3 were first published in this one commit, so the repository
history does not timestamp them separately; it shows only that all of them
existed by that time. Amendment A4 was committed and pushed separately as
\texttt{afd07c0} (full hash
\texttt{\seqsplit{afd07c0b77f9f69c6c2dc873bde0e24a4c43060b}}; commit time 2026-10-30
04:17\,UTC); the design file at that commit has SHA-256
\texttt{\seqsplit{8650efdfa0d282614430c5b0a8d895cf69af0a5f25506e4af5316e20529776ef}}.
A4 appends to the amendment section only; the H100 campaign runs under the
\texttt{ee2e6a7} version of the design, which A4 does not affect. Each
campaign records the SHA-256 of the design document it ran under, and the
analysis resolves it to the oldest commit holding that version: for the H100,
commit \texttt{\stat{meta.design_commit_short}} with \stat{meta.amendments}
amendments; amendments added after it: \stat{meta.later_amendments_list}. The
H100 campaign started at \stat{meta.campaign_start_utc}, after that commit.
% BEGIN-RAN:any
\ifladderran
The other campaigns recorded commit
% BEGIN-RAN:a100
\ifaonehundredran
\texttt{\stat{a100.meta.design_commit_short}} (A100, started
\stat{a100.meta.campaign_start_utc}),
\fi
% END-RAN:a100
% BEGIN-RAN:gh200
\ifghtwohundredran
\texttt{\stat{gh200.meta.design_commit_short}} (GH200, started
\stat{gh200.meta.campaign_start_utc}),
\fi
% END-RAN:gh200
% BEGIN-RAN:b200
\ifbtwohundredran
\texttt{\stat{b200.meta.design_commit_short}} (B200, started
\stat{b200.meta.campaign_start_utc}),
\fi
% END-RAN:b200
each after the commit time it records.
\fi
% END-RAN:any
% Checked 2026-10-30 against the final analysis outputs: every campaign's first
% session started after the commit of the design version it records. An
% external archive (e.g. a Zenodo snapshot of the commits) would be stronger
% evidence than a repository the author controls.

\subsection{Research questions}
\label{sec:rqs}
\begin{description}[leftmargin=*,itemsep=2pt]
  \item[RQ1] How much does the GPU accelerate each framework (Polars GPU versus
    Polars CPU; Spark GPU versus Spark CPU)?
  \item[RQ2] Which engine is fastest per query, suite and scale? Does any GPU
    engine beat Polars CPU?
  \item[RQ3] At which scale from SF1 to SF300, if any, does a GPU engine
    overtake the best CPU engine?
  \item[RQ4] Which queries cannot run on each GPU engine, and why?
  \item[RQ5] (side study) How much does the Polars version (1.44.2 versus
    2.0.0, CPU) change Polars CPU timings?
\end{description}

\subsection{Hosts and the GPU-generation ladder}
\label{sec:hosts}

\begin{table}[t]
\centering
\caption{Hosts (Lambda on-demand instances). The H100 campaign is the primary
analysis; each other host runs its own complete campaign (amendments A2, A3)
and is never pooled with it. Every host runs the engines with 26 threads
(\texttt{--cpus 26}), leaving any further vCPUs idle. Driver and CPU model
are as recorded in each campaign; \hostnotrun{} means the host's campaign did
not run. GPU memory, vCPUs and RAM are the values each host reports
(\texttt{nvidia-smi}, \texttt{lscpu}, \texttt{free}); for a host whose campaign
did not run, the provider's listing, marked~$^{*}$. \stat{hosts.listing_note}
vCPUs: the count the host reports, then the 26 engine threads and the
physical cores they can occupy (\texttt{lscpu} topology, as reported inside
the virtual machine). Driver: version and, in parentheses, the highest CUDA
version it supports;
every host runs the CUDA \stat{hosts.cuda_runtime} runtime.}
\label{tab:hosts}
\scriptsize
\setlength{\tabcolsep}{2.5pt}
\begin{tabular}{@{}llllllp{1.45cm}p{1.9cm}l@{}}
\toprule
Host & Instance type & \shortstack[l]{GPU\\memory} & ISA & \shortstack[l]{vCPUs (threads\\on cores)} & RAM & \shortstack[l]{Driver\\(CUDA)} & CPU model & Campaign \\
\midrule
A100  & \texttt{gpu\_1x\_a100\_sxm4} & \hoststator{a100}{host.gpu_memory_gib}{40$^{*}$}\,GiB & x86\_64 & \hoststator{a100}{host.cpu_layout}{30$^{*}$ (26)} & \hoststator{a100}{host.host_ram_gib}{200$^{*}$}\,GiB
      & \hoststat{a100}{meta.driver_cuda} & \hoststat{a100}{meta.cpu_model} & \stat{hosts.a100.status} \\
H100  & \texttt{gpu\_1x\_h100\_sxm5} & \stat{host.gpu_memory_gib}\,GiB & x86\_64 & \stat{host.cpu_layout} & \stat{host.host_ram_gib}\,GiB
      & \stat{meta.driver_cuda} & \stat{meta.cpu_model} & primary \\
GH200 & \texttt{gpu\_1x\_gh200} & \hoststator{gh200}{host.gpu_memory_gib}{96$^{*}$}\,GiB & aarch64 & \hoststator{gh200}{host.cpu_layout}{64$^{*}$ (26)} & \hoststator{gh200}{host.host_ram_gib}{432$^{*}$}\,GiB
      & \hoststat{gh200}{meta.driver_cuda} & \hoststat{gh200}{meta.cpu_model} & \stat{hosts.gh200.status} \\
B200  & \texttt{gpu\_1x\_b200\_sxm6} & \hoststator{b200}{host.gpu_memory_gib}{180$^{*}$}\,GiB & x86\_64 & \hoststator{b200}{host.cpu_layout}{26$^{*}$ (26)} & \hoststator{b200}{host.host_ram_gib}{360$^{*}$}\,GiB
      & \hoststat{b200}{meta.driver_cuda} & \hoststat{b200}{meta.cpu_model} & \stat{hosts.b200.status} \\
\bottomrule
\end{tabular}
\end{table}

\paragraph{Primary: H100.} All engines ran sequentially on one Lambda
\texttt{gpu\_1x\_h100\_sxm5} instance: one NVIDIA H100 80\,GB SXM5
(\stat{meta.gpu_name}, driver \stat{meta.driver_version}), \stat{host.vcpus} vCPUs
(\stat{meta.cpu_model}), \stat{host.host_ram_gib}\,GiB RAM, and a 2,816\,GiB local volume holding
the data and Spark's local directory. GPU clocks were not locked; they were
recorded (Section~\ref{sec:controls}).

\subsubsection{The ladder (amendments A2, A3)}
\label{sec:ladder-design}
Three further hosts (Table~\ref{tab:hosts}) place the H100 on a ladder of GPU
generations: an A100 40\,GB (Ampere; the GPU of the earlier study), a GH200
(Grace~Hopper: a 96\,GB Hopper GPU joined to an Arm-based Grace CPU by a
coherent NVLink-C2C link~\citep{nvidia_gh200}) and a B200 (Blackwell, 180\,GB). The provider
offered no GB300 or H200 instances, so the B200 is the newest GPU available.
The ladder hosts differ in more than the GPU: CPU model and instruction set,
vCPU count, RAM and the CPU--GPU link all change between rungs.

Each host runs its \emph{own complete campaign} under the unchanged design
(same seed, Latin square, gates and side study), in its own campaign root, with
host identity (instance type, GPU, driver, CPU model, RAM) recorded. Data are
never pooled across hosts; every primary estimate (Section~\ref{sec:stats}) is
computed per host, and the H100 remains the primary analysis. Each host
regenerates its inputs with the same recipe and must pass the same input gate
(Section~\ref{sec:inputs}) and preflight before its round~1.

\paragraph{Same thread count everywhere.} All hosts run the engines with
\texttt{--cpus 26}, so engine thread counts are identical across the ladder;
any further vCPUs (on the A100 and GH200 hosts; Table~\ref{tab:hosts}) are left idle. This favours
comparability across hosts over each host's peak CPU performance
(Section~\ref{sec:threats}). Equal threads are not equal physical cores: where
a host exposes two hardware threads per core, the 26 threads share fewer cores
than on a host without simultaneous multithreading (SMT). On the H100, with
\stat{host.smt_threads_per_core} threads per core, they occupy
\stat{host.cpus_arg_cores_text} physical cores; Table~\ref{tab:hosts} gives
every host's layout. It also differs from the earlier A100 study,
which used 30 threads on the same instance type.

\paragraph{CUDA runtime and drivers.} Every host runs the pinned CUDA
\stat{hosts.cuda_runtime} user-space runtime with the NVIDIA driver of its
provider image (Table~\ref{tab:hosts}); where a host's driver supports only an
older CUDA minor version (here: \stat{hosts.minor_compat_list}), the runtime
runs through CUDA minor-version compatibility, which lets applications built
within one major CUDA release run on an older driver of that release, with
some feature limitations~\citep{nvidia_cuda_compat}; the pinned stack is not
changed.

\paragraph{Conditions for each additional host.} A host's campaign runs only if
the provider has capacity for the instance and the pinned software stack passes
preflight and input verification on it. Software versions are never changed
to make a host run. On the GH200 (aarch64), the same package \emph{versions}
are used where aarch64 builds exist; the platform-specific artifacts (Python
wheels, the plugin JAR's classifier, the JDK build) are recorded with their
SHA-256. If any pinned version has no aarch64 build, the GH200 campaign is not
run. A host that does not run is reported with its reason, below and in
Section~\ref{sec:ladder}; none of these conditions depends on timings. On the
A100 (40\,GB), coverage outcomes from device-memory exhaustion at SF300 are
more likely than on the other hosts; under A1 they are results, not failures.
% BEGIN-NOTRUN:a100
\ifaonehundredran\else
The A100 campaign was not run: \stat{a100.not_run_reason}.
\fi
% END-NOTRUN:a100
% BEGIN-NOTRUN:gh200
\ifghtwohundredran\else
The GH200 campaign was not run: \stat{gh200.not_run_reason}.
\fi
% END-NOTRUN:gh200
% BEGIN-NOTRUN:b200
\ifbtwohundredran\else
The B200 campaign was not run: \stat{b200.not_run_reason}.
\fi
% END-NOTRUN:b200

\subsection{Arms}
\label{sec:arms}

\begin{table}[t]
\centering
\caption{The four arms. No arm was tuned during the campaign; the Spark
configurations are byte-identical to the earlier A100 study.}
\label{tab:arms}
\small
\begin{tabular}{@{}lll@{}}
\toprule
Id & Engine & Version and settings \\
\midrule
\idpolarscpu  & Polars, CPU streaming engine & Polars 1.44.2, Python 3.11, 26 threads \\
\idpolarsgpu  & Polars + cuDF GPU engine & cudf-polars-cu12 26.10.2; \texttt{raise\_on\_fail=True}; \\
            &                         & executor \texttt{streaming} (defaults) \\
\idsparkcpu   & PySpark, \texttt{local[26]} & Spark 3.5.9, Java 17, 96\,GiB driver heap, \\
            &                         & 120 shuffle partitions, AQE on \\
\idsparkgpu   & PySpark + cuDF plugin   & \texttt{rapids-4-spark\_2.12} 26.08.1 (SHA-256 recorded), \\
            &                         & profile \texttt{a100-parity} \\
\bottomrule
\end{tabular}
\end{table}

Table~\ref{tab:arms} lists the arms. Polars runs at version 1.44.2 rather than
the 2.0.0 used on the A100, because no released \texttt{cudf-polars} supports
Polars 2.0: the newest release, 26.10.2, requires Polars ${\ge}1.35,{<}1.45$
and Python ${\ge}3.11$. Both Polars arms share the same Polars build, so the
RQ1 contrast is not confounded by version; RQ5 measures the version effect on
the CPU engine separately.

Both Spark arms share the settings below; the GPU arm adds the plugin JAR
and the \texttt{spark.rapids} and \texttt{spark.plugins} lines:
\begin{flushleft}\small\ttfamily
spark.sql.shuffle.partitions=120\\
spark.sql.adaptive.enabled=true\\
spark.sql.files.maxPartitionBytes=134217728\quad\rmfamily(128\,MiB)\ttfamily\\
spark.sql.session.timeZone=UTC\quad\rmfamily driver heap 96\,GiB\ttfamily\\[2pt]
spark.plugins=com.nvidia.spark.SQLPlugin\\
spark.rapids.sql.enabled=true\\
spark.rapids.sql.concurrentGpuTasks=2\\
spark.rapids.memory.pinnedPool.size=2G\\
spark.rapids.sql.variableFloatAgg.enabled=true\\
spark.rapids.sql.explain=NOT\_ON\_GPU
\end{flushleft}
These are the A100 study's fixed settings, not tuned values. The Polars GPU engine uses
\texttt{cudf-polars}' default streaming-executor options and default memory
resource; its resolved options are recorded per session. Thread counts set
engine worker pools (\texttt{POLARS\_MAX\_THREADS=26}, \texttt{local[26]});
they do not cap total operating-system threads.

\subsection{Workload}
\label{sec:workload}

\begin{table}[t]
\centering
\caption{Queries. Baseline queries are adapted from DuckDB 1.5.6's bundled
TPC-H query text; \q{cx\_q09} and \q{cx\_q13} are derived from TPC-H Q9 and
Q13; the other complex queries are custom. Spark runs the SQL; Polars runs
equivalent native lazy expressions.}
\label{tab:queries}
\small
\begin{tabular}{@{}llp{0.62\linewidth}@{}}
\toprule
Suite & Query & Main work \\
\midrule
\multirow{6}{*}{baseline}
 & \q{q01} & Broad ship-date filter; grouped sums, averages, counts \\
 & \q{q03} & 3-table join, selective predicates, revenue per order, top 10 \\
 & \q{q06} & Selective filter, single revenue sum \\
 & \q{q10} & 4-table join, customer grouping with strings, top 20 \\
 & \q{q12} & 2-table join, conditional counts by ship mode \\
 & \q{q14} & 2-table join, promotional revenue share \\
\midrule
\multirow{6}{*}{complex}
 & \q{cx\_q09} & 6-table join incl.\ composite key, substring filter, nation/year profit \\
 & \q{cx\_q13} & Left outer join, comment pattern filter, nested aggregation \\
 & \q{cx\_semi\_anti} & Correlated \texttt{EXISTS}/\texttt{NOT EXISTS} design (semi/anti joins) \\
 & \q{cx\_window} & \texttt{LAG} and \texttt{ROW\_NUMBER} over two orderings per customer \\
 & \q{cx\_distinct} & 4-table join, three exact distinct counts \\
 & \q{cx\_self\_join} & \texttt{lineitem} self-join with residual inequality \\
\bottomrule
\end{tabular}
\end{table}

Table~\ref{tab:queries} lists the twelve queries. The baseline suite runs at
SF1, 10, 30, 100 and 300 with the A100 study's session lifecycle (no explicit
cleanup between queries). The complex suite (version 2) runs at SF30, 100 and
300 with a \texttt{verified\_gc} barrier after every query action, outside the
timer: Python garbage collection for Polars; Python and JVM garbage collection
for Spark, followed by waiting until the session owns zero shuffle files and
bytes and has no active jobs. A failed barrier aborts the session. Extra sort
keys make top-$N$ ties deterministic.

\paragraph{Type normalization.} Decimal quantity and money columns (the
\texttt{FLOAT\_COLUMNS} set: e.g.\ \texttt{l\_quantity},
\texttt{l\_extendedprice}, \texttt{l\_discount}, \texttt{l\_tax},
\texttt{o\_totalprice}) are cast to Float64/\texttt{DOUBLE} by every engine.
This permits parallel floating-point reductions and does not measure exact
decimal arithmetic.

\subsection{Inputs and verification}
\label{sec:inputs}

Inputs are generated on the host with DuckDB 1.5.6's \texttt{tpch}
extension~\citep{duckdb_tpch_ext} using the A100 study's parameters: SF1--SF30
as a single \texttt{dbgen} call (30 threads, 64\,GB memory limit), SF100 as 20
\texttt{dbgen(children, step)} partitions with 12 workers, and SF300 as 60
partitions with 16 workers. Tables are written as Snappy Parquet. The largest
scale has 1,799,989,091 \texttt{lineitem} rows.

Before round~1, \texttt{verify\_generated\_inputs.py} must report
\texttt{RESULT: OK} for SF1--SF300 (gate result: \stat{meta.input_gate}): row
counts, schemas, generation parameters and order-independent content
fingerprints must match the A100 manifests. Byte identity is not required:
we showed that DuckDB's per-thread output makes the assignment of rows to files
depend on scheduling, so compressed file hashes differ between regenerations
while table contents do not. All engines read the same files and the same
manifest.

\subsection{Sessions, rounds and campaign}
\label{sec:structure}

A \emph{session} is one engine process for one (suite, scale). It preheats the
input files (untimed; residency is not guaranteed), runs one warmup per query
(untimed), and then $R=3$ measured trials per query in a seeded, shuffled
order. The timer covers query construction, optimization, Parquet reading,
execution and result collection, including the device-to-host transfer for
GPU engines. For Polars the timer stops when \texttt{collect} returns a host
DataFrame; conversion with \texttt{to\_dicts()} is outside. For Spark the timer
stops when \texttt{collect()} returns Python \texttt{Row} objects; conversion
with \texttt{asDict()} is outside (Section~\ref{sec:threats}).

A \emph{round} runs every engine on every (suite, scale) once: $4 \times 8 = 32$
sessions. The \emph{campaign} is four rounds (128 sessions); the round is the
block for analysis. The planned measured trials number
$4 \text{ rounds} \times (6 \times 5 + 6 \times 3) \text{ cases} \times 4
\text{ engines} \times 3 = 2{,}304$ executions, of which those for the two
pre-declared unsupported \polarsgpu{} queries (Section~\ref{sec:coverage-design})
are recorded as errors.

\subsection{Order design}
\label{sec:latin}

\begin{table}[t]
\centering
\caption{Engine order: a Williams balanced Latin square for $n=4$, with
letters assigned by \texttt{random.Random(20261029).shuffle} (A~=~\idsparkcpu,
B~=~\idpolarsgpu, C~=~\idpolarscpu, D~=~\idsparkgpu). Each engine takes each position
once; each ordered pair of engines is adjacent exactly once.}
\label{tab:latin}
\small
\begin{tabular}{@{}cclc@{}}
\toprule
Round & Sequence & Engine order within every (suite, scale) & First suite \\
\midrule
1 & A B D C & \sparkcpu{} $\to$ \polarsgpu{} $\to$ \sparkgpu{} $\to$ \polarscpu{} & baseline \\
2 & B C A D & \polarsgpu{} $\to$ \polarscpu{} $\to$ \sparkcpu{} $\to$ \sparkgpu{} & complex \\
3 & C D B A & \polarscpu{} $\to$ \sparkgpu{} $\to$ \polarsgpu{} $\to$ \sparkcpu{} & baseline \\
4 & D A C B & \sparkgpu{} $\to$ \sparkcpu{} $\to$ \polarscpu{} $\to$ \polarsgpu{} & complex \\
\bottomrule
\end{tabular}
\end{table}

Engine order follows the Williams design~\citep{williams1949latin} in
Table~\ref{tab:latin}, seeded with 20261029. The twelve adjacent ordered pairs
(AB, BD, DC; BC, CA, AD; CD, DB, BA; DA, AC, CB) are all distinct, so
first-order carryover (for example, page-cache or thermal state left by the
previous engine) is balanced across engines. Suite order alternates between
rounds; scales run in ascending order within a suite; within-session query
order uses the worker's fixed seed.

\subsection{Controls}
\label{sec:controls}
\begin{itemize}[leftmargin=*,itemsep=1pt]
  \item No other workload ran on the host during rounds; data generation was
    finished and verified first.
  \item Before each session: no GPU compute processes, input files preheated,
    previous session's processes exited.
  \item Telemetry: GPU utilization, memory, power, temperature and SM/memory
    clocks every 200\,ms; host CPU, memory, swap and disk every 1\,s.
  \item \texttt{--cpus 26} for all engines; same data root and manifest.
  \item After every round, results were copied off the host and checksummed.
\end{itemize}

\subsection{Integrity gates}
\label{sec:gates}

Each round must pass all of the following, or it is repeated once and the
repeat is reported; a second failure stops the campaign. None depends on
timings.
\begin{enumerate}[leftmargin=*,itemsep=1pt]
  \item Every session exits; no engine timeout; no unclassified error.
  \item Every output agrees with \polarscpu{} for the same query: relative
    tolerance $10^{-8}$, absolute $10^{-6}$, integers exact.
  \item \polarsgpu{}: the strict GPU session guard passes, and every query
    failure is a classified coverage outcome (below). The guard confirms the resolved
    executor and \texttt{raise\_on\_fail=True}, and runs a canary query with a
    Python user-defined function, which must raise; if it does not, the
    engine could be falling back to the CPU silently.
  \item \sparkgpu{}: final plans contain GPU scan and GPU compute operators and
    no CPU core operators (complex rule), or GPU scan and aggregation
    (baseline rule).
  \item Complex suite: every \texttt{verified\_gc} barrier passes.
\end{enumerate}
\paragraph{Coverage outcomes versus integrity failures (amendment A1).}
Failures are classified, and only integrity failures fail a round. A
\emph{coverage outcome} is a GPU engine rejecting or failing to complete a
query or session because of (a) an unsupported operation (for example a
\texttt{cudf-polars} \texttt{NotImplementedError} under
\texttt{raise\_on\_fail=True}) or (b) resource exhaustion (device or host
out of memory, an RMM allocation failure, or Spark GPU out-of-memory after
retries). The same rule covers a CPU engine that exhausts host resources.
Coverage outcomes are recorded with their error class and scale, reported
under RQ4, and excluded from contrasts involving that engine. An
\emph{integrity failure} is an output disagreement, a silent CPU fallback or
CPU core operators in a GPU plan, a failed \texttt{verified\_gc} barrier, an
engine timeout, an unclassified error, or a crash not attributable to (a) or
(b). If a coverage outcome occurs in some rounds but not others for the same
engine, suite, scale and query, the cell is \emph{unstable coverage} and is
excluded from contrasts in all rounds. The classification rules (error
type and message patterns) are fixed in the preflight tool, and any error
that matches no rule counts as an integrity failure.

A post-round-1 review checked integrity only. No stopping or modification was
based on observed speeds.

\subsection{Pre-declared GPU coverage gaps}
\label{sec:coverage-design}

Preflight on 2026-10-26 showed that \polarsgpu{} rejects \q{cx\_q13} (the
inline regular-expression flag \texttt{(?s)} is unsupported) and
\q{cx\_window} (\texttt{cum\_count().over(...)} is unsupported) under
\texttt{raise\_on\_fail=True}. Under A1 they are coverage outcomes of type
(a): reported as not supported on the GPU (RQ4) and excluded from ratios
involving \polarsgpu{}. Queries are never rewritten for the GPU within this
campaign.

\subsection{Statistical analysis plan}
\label{sec:stats}

Let $t_{e,q,r}$ be the median of the three measured trials of engine $e$ on
query $q$ in round $r$ (for a fixed suite and scale). Each contrast compares a
reference engine $a$ with an engine $b$ as the speedup ratio
$\rho_{q,r} = t_{a,q,r} / t_{b,q,r}$, so $\rho > 1$ means $b$ is faster:
\begin{itemize}[leftmargin=*,itemsep=1pt]
  \item \textbf{C1} \polarscpu{}\,/\,\polarsgpu{} (RQ1: Polars GPU speedup);
  \item \textbf{C2} \sparkcpu{}\,/\,\sparkgpu{} (RQ1: Spark GPU speedup);
  \item \textbf{C3} \sparkgpu{}\,/\,\polarsgpu{} (RQ2: GPU engine head-to-head);
  \item \textbf{C4} \polarscpu{}\,/\,best GPU engine, the faster of the
    GPU engines that ran the query in that round (RQ2, RQ3). A stricter
    variant, C4$_\text{strict}$, limits C4 to queries that both GPU engines
    ran; it is reported in the artifact.
\end{itemize}
Per query, let $\ell_{q,r} = \log \rho_{q,r}$. The estimate is
$\exp(\bar\ell_q)$, the geometric mean over the four rounds; the 95\%
interval is $\exp\!\bigl(\bar\ell_q \pm t_{0.975,3}\, s_q / \sqrt{4}\bigr)$
with $t_{0.975,3} \approx 3.182$ and $s_q$ the standard deviation of
$\ell_{q,r}$ across rounds. We also report the across-round minimum and
maximum. Per (suite, scale), the per-round suite statistic is the mean of
$\ell_{q,r}$ over the queries supported by both engines, and the same
$t$-interval is taken across rounds (the round is the block). Geometric means
of ratios follow~\citet{fleming1986geomean}. For comparison with the A100 study
we also report legacy sums of per-query medians, never mixed with geometric
means.

Secondary, descriptive analyses: within-session (across trials) and
between-round coefficients of variation per engine; mean log-time by session
position (1--4) and by predecessor engine (order and carryover check); peak GPU
memory; process wall time and startup; cleanup cost; an energy proxy (mean
board power $\times$ time); and a comparison of the GPU Spark engine between
the earlier A100 study and the H100 campaign, labelled as a comparison of
non-equivalent hosts. CPU engines are not compared with the earlier study
(30 versus 26 threads); the A100 ladder campaign below is the like-for-like
comparison.

\paragraph{Cross-host analysis (amendments A2, A3, secondary).} For each
additional host whose campaign runs, all primary estimates above are computed
from that host's campaign alone. Cross-host comparison is descriptive and uses
the H100 as the reference: for each host $h$, suite, scale, query and engine,
the ratio of the per-host geometric-mean round medians,
$t_{h}/t_{\mathrm{H100}}$ (below 1: faster on host $h$), with an interval from
the two hosts' round distributions (Welch-type, on the log scale); the suite
summary is the geometric mean over queries supported on both hosts. For GPU
engines, a cross-host difference confounds the GPU with the CPU, memory, CPU--GPU
link and driver. CPU engines are compared across hosts only with the caveat
that the CPU model, instruction set and RAM differ; all are recorded.

\textbf{RQ3 crossover} is the smallest scale, per suite, at which the C4
interval lies entirely above 1, if any. \textbf{Multiplicity:} four contrasts
$\times$ suites $\times$ scales are reported as estimates with intervals; no
hypothesis tests or $p$-values are used for claims. With four rounds and
$\text{df}=3$, intervals are expected to be wide.

\subsection{Side study: Polars version bridge (RQ5)}
\label{sec:bridge-design}
After the main rounds, on the same host, \polarscpu{} 1.44.2 (Python 3.11) is
compared with Polars 2.0.0 (Python 3.10, the A100 environment) on both suites
and all scales, in four rounds ordered AB, BA, BA, AB, using the same session
protocol. The estimate is the geometric-mean ratio
$t_{2.0.0} / t_{1.44.2}$ per suite and scale with the same $t$-interval.

\subsection{Scope, deviations and notice}
\label{sec:scope}
\tpcdisclaimer{} Deviations from the TPC-H specification include: a selected
subset of six adapted queries with fixed substitution parameters, plus two
TPC-H-derived and four custom complex queries; data from DuckDB's generator,
which differs slightly from the specification~\citep{duckdb_tpch_ext};
Float64 instead of decimal arithmetic; extra sort keys; Parquet files rather
than a loaded database; no refresh functions, power or throughput tests, or
TPC metrics; and our own timing boundary and summary statistics. No result in
this paper may be compared with official TPC-H results. Use of the benchmark
name and the TPC software is subject to the TPC
policies~\citep{tpcpolicies620,tpcfairuse2019} and the TPC
EULA~\citep{tpceula22}.
% TODO(author): confirm TPC written-permission status for benchmark-name use
% (Policies 6.20 sec. 8.1.4) and the final name of the derived workload
% (\workloadname in macros.tex). See paper/README.md.

% ======================================================================
\section{Results}
\label{sec:results}

% Every result sentence is drawn from \stat{} keys; claim strength follows the
% interval verdicts (above 1, below 1, includes 1).

All results in this section come from the primary H100 campaign, except
Section~\ref{sec:ladder}, which reports the other hosts.

\paragraph{Campaign integrity.} \stat{meta.rounds_completed} rounds completed
(\stat{meta.rounds_excluded_count} excluded, \stat{meta.rounds_repeated} repeated under the gate rule) with
\stat{meta.sessions_completed} sessions and \stat{meta.measured_trials}
measured trials. All \stat{meta.outputs_compared} saved outputs were compared
with \polarscpu{}; \stat{meta.outputs_failed} failed (largest relative
difference \stat{meta.max_rel_diff}). \stat{meta.spark_gpu_plans_checked}
\sparkgpu{} final plans passed the plan rule, and
\stat{meta.cleanup_barriers_passed} \texttt{verified\_gc} barriers passed. The
largest relative difference from \polarscpu{} was
\stat{agreement_max.polars_gpu.max_rel_diff} for \polarsgpu{}
(\stat{agreement_max.polars_gpu.where}) and
\stat{agreement_max.spark_gpu.max_rel_diff} for \sparkgpu{}
(\stat{agreement_max.spark_gpu.where}), against a relative tolerance of
$10^{-8}$.
Campaign window: \stat{meta.campaign_start_utc} to \stat{meta.campaign_end_utc}.

\begin{figure}[t]
\centering
\studyfig{f10_headline_times}{\linewidth}
\caption{Per-query median time for each engine at SF300 (geometric mean over
four rounds of per-round medians; one marker per engine). Bars span the fastest
to slowest round and are hidden behind the marker where the rounds agree
closely. A query with no timing on an engine appears as a hollow marker of
that engine in the hatched ``no timing'' strip below the axis, labelled with
the reason (here: not supported). Log scale; lower is better. \chartnote}
\label{fig:f10}
\end{figure}

\begin{table}[t]
\centering
\caption{Suite-level geometric-mean speedup ratios with 95\% CIs ($t$,
df~=~3, across four rounds). $>1$: the second-named engine is faster. C1 and
C3 exclude \q{cx\_q13} and \q{cx\_window} in the complex suite. \chartnote}
\label{tab:contrasts}
\scriptsize
\setlength{\tabcolsep}{4pt}
\begin{tabular}{@{}llcccc@{}}
\toprule
Suite & Scale & \shortstack{C1\\\polarscpu/\\\polarsgpu} & \shortstack{C2\\\sparkcpu/\\\sparkgpu} & \shortstack{C3\\\sparkgpu/\\\polarsgpu} & \shortstack{C4\\\polarscpu/\\best GPU} \\
\midrule
baseline & SF1   & \ratioci{c1.baseline.sf1}   & \ratioci{c2.baseline.sf1}   & \ratioci{c3.baseline.sf1}   & \ratioci{c4.baseline.sf1} \\
         & SF10  & \ratioci{c1.baseline.sf10}  & \ratioci{c2.baseline.sf10}  & \ratioci{c3.baseline.sf10}  & \ratioci{c4.baseline.sf10} \\
         & SF30  & \ratioci{c1.baseline.sf30}  & \ratioci{c2.baseline.sf30}  & \ratioci{c3.baseline.sf30}  & \ratioci{c4.baseline.sf30} \\
         & SF100 & \ratioci{c1.baseline.sf100} & \ratioci{c2.baseline.sf100} & \ratioci{c3.baseline.sf100} & \ratioci{c4.baseline.sf100} \\
         & SF300 & \ratioci{c1.baseline.sf300} & \ratioci{c2.baseline.sf300} & \ratioci{c3.baseline.sf300} & \ratioci{c4.baseline.sf300} \\
\midrule
complex  & SF30  & \ratioci{c1.complex.sf30}   & \ratioci{c2.complex.sf30}   & \ratioci{c3.complex.sf30}   & \ratioci{c4.complex.sf30} \\
         & SF100 & \ratioci{c1.complex.sf100}  & \ratioci{c2.complex.sf100}  & \ratioci{c3.complex.sf100}  & \ratioci{c4.complex.sf100} \\
         & SF300 & \ratioci{c1.complex.sf300}  & \ratioci{c2.complex.sf300}  & \ratioci{c3.complex.sf300}  & \ratioci{c4.complex.sf300} \\
\bottomrule
\end{tabular}
\end{table}

\subsection{RQ1: GPU speedup within each framework}
\label{sec:rq1}

\begin{figure}[t]
\centering
\studyfig{f01_forest_sf300}{\linewidth}
\caption{Per-query speedup ratios at SF300 (geometric mean over four rounds)
with 95\% CIs, by contrast (C1--C4) and suite; the other scales are in
Appendix~\ref{app:forest}. Vertical line at 1: no difference;
right of 1: the second-named engine is faster. Log scale. Intervals have
df~=~3 and are wide by design. \chartnote}
\label{fig:f1}
\end{figure}

Table~\ref{tab:contrasts} and Figure~\ref{fig:f1} give the C1 and C2 estimates.
On the baseline suite at SF300, the Polars GPU engine's geometric-mean speedup
over Polars CPU was \ratiotext{c1.baseline.sf300}; the Spark plugin's speedup
over CPU Spark was \ratiotext{c2.baseline.sf300}. On the complex suite at
SF300 the corresponding values were \ratiotext{c1.complex.sf300} (over
\stat{c1.complex.sf300.n_queries} supported queries) and
\ratiotext{c2.complex.sf300}. At SF1, C1 was \ratiotext{c1.baseline.sf1} and C2
was \ratiotext{c2.baseline.sf1}.

On the baseline suite, C1 \stat{c1.baseline.summary.trend}; its interval lay
below~1 at \stat{c1.baseline.summary.below} (Polars CPU faster) and above~1 at
\stat{c1.baseline.summary.above}. On the complex suite C1 ranged from
\stat{c1.complex.summary.min} to \stat{c1.complex.summary.max} and its interval
lay above~1 at \stat{c1.complex.summary.above}; it
\stat{c1.complex.summary.trend}. C2 behaved similarly: on the baseline suite
its interval was \stat{c2.baseline.summary.pattern}; on the complex suite it
ranged from
\stat{c2.complex.summary.min} to \stat{c2.complex.summary.max}, above~1 at
\stat{c2.complex.summary.above}. Comparing suites,
\stat{c12.complex_larger_text}. Per query (Figure~\ref{fig:f1}), the smallest
C1 at SF300 was the selective scan-and-sum \q{q06}
(\ratiotext{c1.baseline.sf300.q.q06}; interval
\stat{c1.baseline.sf300.q.q06.verdict}), against
\ratiotext{c1.baseline.sf300.q.q01} for the grouped aggregation \q{q01} and
\ratiotext{c1.baseline.sf300.q.q03} for the three-table join \q{q03}.

\paragraph{A version-locked CPU baseline.} C1 and C4 use Polars 1.44.2 on the
CPU, the version the GPU engine requires, not the current release. In the
side study (Section~\ref{sec:rq5}), Polars 2.0.0 on the CPU was
\stat{bridge_summary.both}; the most extreme query was
\q{\stat{bridge_summary.complex.extreme_query}}, from
\stat{bridge_summary.complex.extreme.min} to
\stat{bridge_summary.complex.extreme.max} across scales. C1 is therefore the
GPU speedup over the CPU engine of the same Polars build, not over the
fastest Polars CPU release available. A post-hoc comparison against Polars
2.0.0 is reported separately as exploratory
(Section~\ref{sec:exploratory}).
% The direction wording is generated from the intervals (statskeys._bridge_summary).

\subsection{RQ2: Fastest engine per query, suite and scale}
\label{sec:rq2}

\begin{figure}[t]
\centering
\begin{subfigure}[t]{\linewidth}
  \centering\studyfig{f03_heatmap_baseline}{\linewidth}
  \caption{Baseline suite.}
\end{subfigure}\\[4pt]
\begin{subfigure}[t]{\linewidth}
  \centering\studyfig{f03_heatmap_complex}{\linewidth}
  \caption{Complex suite.}
\end{subfigure}
\caption{Median time (seconds) for each query, scale and engine, with the
fastest engine marked per row. Grey: no timing (coverage outcome on that GPU
engine). \chartnote}
\label{fig:f3}
\end{figure}

Across the \stat{fastest.total.baseline.count} baseline query--scale cells,
the fastest engine (lowest geometric-mean median over rounds) was \polarscpu{}
in \stat{fastest.polars.baseline.count}, \polarsgpu{} in
\stat{fastest.polars_gpu.baseline.count}, \sparkcpu{} in
\stat{fastest.spark_cpu.baseline.count} and \sparkgpu{} in
\stat{fastest.spark_gpu.baseline.count} (Figure~\ref{fig:f3}). Across the
\stat{fastest.total.complex.count} complex cells the counts were
\stat{fastest.polars.complex.count}, \stat{fastest.polars_gpu.complex.count},
\stat{fastest.spark_cpu.complex.count} and
\stat{fastest.spark_gpu.complex.count}. The two GPU engines compared
head-to-head (C3) at \ratiotext{c3.baseline.sf300} on the baseline suite and
\ratiotext{c3.complex.sf300} on the complex suite at SF300. Against Polars
CPU, the best GPU engine per query (C4) reached \ratiotext{c4.baseline.sf300}
and \ratiotext{c4.complex.sf300}. A ``fastest'' count does not imply a
resolved difference; Figure~\ref{fig:f1} shows which per-query intervals
exclude~1. \polarsgpu{} won \stat{fastest.polars_gpu.complex.of_supported}
complex cells it could run; the complex cells it did not win were won by
\sparkgpu{} (\stat{fastest.spark_gpu.complex.won_text}) and \polarscpu{}
(\stat{fastest.polars.complex.won_text}). On the baseline suite, \polarscpu{}
won its cells at \stat{fastest.polars.baseline.won_text}, and \polarsgpu{} the
others. \sparkcpu{} won no cell on either suite. Between the two GPU engines,
C3 ranged from \stat{c3.baseline.summary.min} to \stat{c3.baseline.summary.max}
on the baseline suite (it \stat{c3.baseline.summary.trend}) and from
\stat{c3.complex.summary.min} to \stat{c3.complex.summary.max} on the complex
suite, with intervals above~1 \stat{c3.both.above_text}.

Table~\ref{tab:sf300} reports per-query medians at SF300, and
Table~\ref{tab:legacy} the legacy sums of medians.

\begin{table}[t]
\centering
\caption{SF300 per-query medians in seconds (geometric mean over four
rounds of per-round medians). ``n/s'': not supported on that GPU engine;
``unstable'': unstable coverage (a coverage outcome in some rounds only),
excluded under A1.
\chartnote}
\label{tab:sf300}
\small
\begin{tabular}{@{}lrrrr@{}}
\toprule
Query & \polarscpu & \polarsgpu & \sparkcpu & \sparkgpu \\
\midrule
\q{q01} & \stat{median.polars.baseline.sf300.q01} & \stat{median.polars_gpu.baseline.sf300.q01} & \stat{median.spark_cpu.baseline.sf300.q01} & \stat{median.spark_gpu.baseline.sf300.q01} \\
\q{q03} & \stat{median.polars.baseline.sf300.q03} & \stat{median.polars_gpu.baseline.sf300.q03} & \stat{median.spark_cpu.baseline.sf300.q03} & \stat{median.spark_gpu.baseline.sf300.q03} \\
\q{q06} & \stat{median.polars.baseline.sf300.q06} & \stat{median.polars_gpu.baseline.sf300.q06} & \stat{median.spark_cpu.baseline.sf300.q06} & \stat{median.spark_gpu.baseline.sf300.q06} \\
\q{q10} & \stat{median.polars.baseline.sf300.q10} & \stat{median.polars_gpu.baseline.sf300.q10} & \stat{median.spark_cpu.baseline.sf300.q10} & \stat{median.spark_gpu.baseline.sf300.q10} \\
\q{q12} & \stat{median.polars.baseline.sf300.q12} & \stat{median.polars_gpu.baseline.sf300.q12} & \stat{median.spark_cpu.baseline.sf300.q12} & \stat{median.spark_gpu.baseline.sf300.q12} \\
\q{q14} & \stat{median.polars.baseline.sf300.q14} & \stat{median.polars_gpu.baseline.sf300.q14} & \stat{median.spark_cpu.baseline.sf300.q14} & \stat{median.spark_gpu.baseline.sf300.q14} \\
\midrule
\q{cx\_q09} & \stat{median.polars.complex.sf300.cx_q09} & \stat{median.polars_gpu.complex.sf300.cx_q09} & \stat{median.spark_cpu.complex.sf300.cx_q09} & \stat{median.spark_gpu.complex.sf300.cx_q09} \\
\q{cx\_q13} & \stat{median.polars.complex.sf300.cx_q13} & \stat{median.polars_gpu.complex.sf300.cx_q13} & \stat{median.spark_cpu.complex.sf300.cx_q13} & \stat{median.spark_gpu.complex.sf300.cx_q13} \\
\q{cx\_semi\_anti} & \stat{median.polars.complex.sf300.cx_semi_anti} & \stat{median.polars_gpu.complex.sf300.cx_semi_anti} & \stat{median.spark_cpu.complex.sf300.cx_semi_anti} & \stat{median.spark_gpu.complex.sf300.cx_semi_anti} \\
\q{cx\_window} & \stat{median.polars.complex.sf300.cx_window} & \stat{median.polars_gpu.complex.sf300.cx_window} & \stat{median.spark_cpu.complex.sf300.cx_window} & \stat{median.spark_gpu.complex.sf300.cx_window} \\
\q{cx\_distinct} & \stat{median.polars.complex.sf300.cx_distinct} & \stat{median.polars_gpu.complex.sf300.cx_distinct} & \stat{median.spark_cpu.complex.sf300.cx_distinct} & \stat{median.spark_gpu.complex.sf300.cx_distinct} \\
\q{cx\_self\_join} & \stat{median.polars.complex.sf300.cx_self_join} & \stat{median.polars_gpu.complex.sf300.cx_self_join} & \stat{median.spark_cpu.complex.sf300.cx_self_join} & \stat{median.spark_gpu.complex.sf300.cx_self_join} \\
\bottomrule
\end{tabular}
\end{table}

\begin{table}[t]
\centering
\caption{Legacy score: sum of per-query medians (seconds), averaged over
rounds, for comparability with the A100 study's reporting. Not a once-through
suite time and not comparable with the geometric means above. A total is
reported only when every query has a valid median in every round, so the
\polarsgpu{} complex totals are n/a (two unsupported queries). \chartnote}
\label{tab:legacy}
\small
\begin{tabular}{@{}llrrrr@{}}
\toprule
Suite & Scale & \polarscpu & \polarsgpu & \sparkcpu & \sparkgpu \\
\midrule
baseline & SF1   & \stat{legacy.sum_medians.polars.baseline.sf1}   & \stat{legacy.sum_medians.polars_gpu.baseline.sf1}   & \stat{legacy.sum_medians.spark_cpu.baseline.sf1}   & \stat{legacy.sum_medians.spark_gpu.baseline.sf1} \\
         & SF10  & \stat{legacy.sum_medians.polars.baseline.sf10}  & \stat{legacy.sum_medians.polars_gpu.baseline.sf10}  & \stat{legacy.sum_medians.spark_cpu.baseline.sf10}  & \stat{legacy.sum_medians.spark_gpu.baseline.sf10} \\
         & SF30  & \stat{legacy.sum_medians.polars.baseline.sf30}  & \stat{legacy.sum_medians.polars_gpu.baseline.sf30}  & \stat{legacy.sum_medians.spark_cpu.baseline.sf30}  & \stat{legacy.sum_medians.spark_gpu.baseline.sf30} \\
         & SF100 & \stat{legacy.sum_medians.polars.baseline.sf100} & \stat{legacy.sum_medians.polars_gpu.baseline.sf100} & \stat{legacy.sum_medians.spark_cpu.baseline.sf100} & \stat{legacy.sum_medians.spark_gpu.baseline.sf100} \\
         & SF300 & \stat{legacy.sum_medians.polars.baseline.sf300} & \stat{legacy.sum_medians.polars_gpu.baseline.sf300} & \stat{legacy.sum_medians.spark_cpu.baseline.sf300} & \stat{legacy.sum_medians.spark_gpu.baseline.sf300} \\
\midrule
complex  & SF30  & \stat{legacy.sum_medians.polars.complex.sf30}   & \stat{legacy.sum_medians.polars_gpu.complex.sf30}   & \stat{legacy.sum_medians.spark_cpu.complex.sf30}   & \stat{legacy.sum_medians.spark_gpu.complex.sf30} \\
         & SF100 & \stat{legacy.sum_medians.polars.complex.sf100}  & \stat{legacy.sum_medians.polars_gpu.complex.sf100}  & \stat{legacy.sum_medians.spark_cpu.complex.sf100}  & \stat{legacy.sum_medians.spark_gpu.complex.sf100} \\
         & SF300 & \stat{legacy.sum_medians.polars.complex.sf300}  & \stat{legacy.sum_medians.polars_gpu.complex.sf300}  & \stat{legacy.sum_medians.spark_cpu.complex.sf300}  & \stat{legacy.sum_medians.spark_gpu.complex.sf300} \\
\bottomrule
\end{tabular}
\end{table}

\subsection{RQ3: Scale at which a GPU engine overtakes the best CPU engine}
\label{sec:rq3}

\begin{figure}[t]
\centering
\studyfig{f02_suite_speedup_vs_scale}{\linewidth}
\caption{Suite geometric-mean speedup versus scale factor for C1--C4, with 95\%
CI bands (df~=~3). Lines connect measured scales only; they are not a fitted
model and do not extrapolate. \chartnote}
\label{fig:f2}
\end{figure}

By the pre-registered rule (the smallest scale at which the C4 interval lies
entirely above~1), the baseline crossover scale was
\stat{crossover.baseline.c4.scale} and the complex crossover scale was
\stat{crossover.complex.c4.scale} (Figure~\ref{fig:f2}).
Below the crossover, C4's interval lay below~1 at
\stat{c4.baseline.summary.below} on the baseline suite, so there Polars CPU was
faster than either GPU engine. On the complex suite the interval lay above~1 at
\stat{c4.complex.summary.above}. A crossover reported at a suite's smallest
measured scale (SF1 for the baseline suite, SF30 for the complex suite) only
says that the GPU was already ahead there; a crossover between two measured
scales is not localized more precisely than that interval. Per-query
crossovers are visible in Figure~\ref{fig:f1}.

\subsection{RQ4: GPU coverage}
\label{sec:rq4}

\begin{figure}[t]
\centering
\studyfig{f04_gpu_coverage}{0.8\linewidth}
\caption{GPU coverage by query, scale and GPU engine (amendment A1
classes): ran on the GPU with all gates passed; not supported (unsupported
operation); resource exhaustion (for example out of device memory); other
error (not a coverage outcome; an integrity failure under A1); or partial
(hatched: unstable coverage, a coverage outcome in some rounds only).}
\label{fig:f4}
\end{figure}

\polarsgpu{} ran \stat{coverage.polars_gpu.baseline.supported} of
\stat{coverage.polars_gpu.baseline.total} baseline and
\stat{coverage.polars_gpu.complex.supported} of
\stat{coverage.polars_gpu.complex.total} complex query--scale cells on the GPU
(Figure~\ref{fig:f4}). The unsupported queries were
\stat{coverage.polars_gpu.unsupported_list}; the number of failures not
pre-declared in Section~\ref{sec:coverage-design} was
\stat{coverage.polars_gpu.unexpected_failures}. \sparkgpu{} ran
\stat{coverage.spark_gpu.baseline.supported} of
\stat{coverage.spark_gpu.baseline.total} and
\stat{coverage.spark_gpu.complex.supported} of
\stat{coverage.spark_gpu.complex.total} cells with plans meeting the GPU plan
rule. Both pre-declared \polarsgpu{} gaps are expression-level: the
\texttt{(?s)} inline regex flag in \q{cx\_q13}, and
\texttt{cum\_count().over(...)} in \q{cx\_window}, whose
\texttt{GroupedWindow} translation in \texttt{cudf-polars} 26.10.2 accepts only
a fixed set of window functions. Resource exhaustion accounted for
\stat{coverage.polars_gpu.baseline.resource_exhaustion} baseline and
\stat{coverage.polars_gpu.complex.resource_exhaustion} complex \polarsgpu{}
cells and \stat{coverage.spark_gpu.baseline.resource_exhaustion} and
\stat{coverage.spark_gpu.complex.resource_exhaustion} \sparkgpu{} cells;
unstable coverage, \stat{coverage.polars_gpu.baseline.unstable},
\stat{coverage.polars_gpu.complex.unstable},
\stat{coverage.spark_gpu.baseline.unstable} and
\stat{coverage.spark_gpu.complex.unstable} cells respectively. Across the
\stat{ladder.hosts_count} hosts of the ladder, \stat{ladder.resource_exhaustion_total}
GPU-engine cells were lost to resource exhaustion and
\stat{ladder.unexpected_failures_total} \polarsgpu{} failures were not
pre-declared (Table~\ref{tab:ladder-coverage}).

\begin{figure}[t]
\centering
\studyfig{f07_gpu_memory}{0.8\linewidth}
\caption{GPU memory per query and GPU engine (maximum over rounds); the two
engines' values measure different things. \polarsgpu{}: the worker process's
device memory sampled every 200\,ms during the timed action; cudf-polars'
memory pool keeps freed blocks, so within a session where the sample never
falls it is the pool's high-water mark, an upper bound on what a query needs
rather than its working set (see text for how many sessions). \sparkgpu{}: the
task-reported maximum, the largest device allocation that a single task
reported (RAPIDS task metric) among the tasks of the action; not a total over
concurrent tasks. The plugin's pool reservation is drawn as a line only for a
host where the sampled value was constant in every session. Sampling can miss
short peaks.}
\label{fig:f7}
\end{figure}

The two GPU engines' device memory is measured differently
(Figure~\ref{fig:f7}). For \polarsgpu{}, the sampled device memory of the
worker process at SF300 reached
\stat{peak_gpu_mem_gib.polars_gpu.baseline.sf300.max}\,GiB (baseline) and
\stat{peak_gpu_mem_gib.polars_gpu.complex.sf300.max}\,GiB (complex). cudf-polars
allocates from a memory pool that keeps freed blocks; on the H100,
\stat{gpu_mem_note.polars_gpu}. For \sparkgpu{}, the
sampled value (\stat{pool_gpu_mem_gib.spark_gpu.complex.sf300.max}\,GiB on the
complex suite at SF300) does not follow the action's own use: on the H100,
\stat{gpu_mem_note.spark_gpu}. We therefore report the task-reported maximum instead: the
largest device allocation that a single task reported among the tasks of the
timed action, \stat{task_gpu_mem_gib.spark_gpu.baseline.sf300.max}\,GiB
(baseline) and \stat{task_gpu_mem_gib.spark_gpu.complex.sf300.max}\,GiB
(complex) at SF300, available for
\stat{secondary.gpu_memory.task_reported_cells.spark_gpu.with_value} of
\stat{secondary.gpu_memory.task_reported_cells.spark_gpu.cells} query--scale
cells. It is not a total over concurrently running tasks, so neither engine's
value is the minimum memory a query needs.

\subsection{RQ5: Polars version bridge}
\label{sec:rq5}

\begin{figure}[t]
\centering
\studyfig{f08_polars_version_bridge}{0.8\linewidth}
\caption{Polars CPU version bridge: geometric-mean ratio
$t_{2.0.0}/t_{1.44.2}$ per suite and scale with 95\% CIs (four rounds, orders
AB, BA, BA, AB). Above 1: 1.44.2 is faster. \chartnote}
\label{fig:f8}
\end{figure}

\begin{table}[t]
\centering
\caption{Polars version bridge, $t_{2.0.0}/t_{1.44.2}$, geometric mean with
95\% CI.}
\label{tab:bridge}
\small
\begin{tabular}{@{}lccccc@{}}
\toprule
Suite & SF1 & SF10 & SF30 & SF100 & SF300 \\
\midrule
baseline & \ratioci{bridge.baseline.sf1} & \ratioci{bridge.baseline.sf10} & \ratioci{bridge.baseline.sf30} & \ratioci{bridge.baseline.sf100} & \ratioci{bridge.baseline.sf300} \\
complex  & --- & --- & \ratioci{bridge.complex.sf30} & \ratioci{bridge.complex.sf100} & \ratioci{bridge.complex.sf300} \\
\bottomrule
\end{tabular}
\end{table}

The version ratio at SF300 was \ratiotext{bridge.baseline.sf300} on the
baseline suite and \ratiotext{bridge.complex.sf300} on the complex suite
(Table~\ref{tab:bridge}, Figure~\ref{fig:f8}). Polars~2.0.0 was
\stat{bridge_summary.both}. The effect is
not small next to the engine contrasts: it bears directly on C1 and C4, whose
CPU side is the older version (Section~\ref{sec:rq1}), and on any comparison
with the earlier A100 study, which used Polars \stat{a100_h100.study.polars_version}.
Each side-study estimate rests on
\stat{bridge_summary.baseline.n_rounds} accepted rounds (four were planned,
orders AB, BA, BA, AB).
% Four side-study rounds were accepted on the H100 (checked 2026-10-30).

\subsection{Secondary analyses}
\label{sec:secondary}

\begin{figure}[t]
\centering
\begin{subfigure}[t]{0.49\linewidth}
  \centering\studyfig{f05_variability}{\linewidth}
  \caption{Coefficient of variation within sessions and between rounds.}
  \label{fig:f5}
\end{subfigure}\hfill
\begin{subfigure}[t]{0.49\linewidth}
  \centering\studyfig{f06_order_effects}{\linewidth}
  \caption{Mean log-time deviation by session position and predecessor.}
  \label{fig:f6}
\end{subfigure}
\caption{Variability and order checks (descriptive). \chartnote}
\label{fig:f56}
\end{figure}

\paragraph{Variability.} The median within-session CV was
\stat{cv.within.polars.median}\% (\polarscpu{}),
\stat{cv.within.polars_gpu.median}\% (\polarsgpu{}),
\stat{cv.within.spark_cpu.median}\% (\sparkcpu{}) and
\stat{cv.within.spark_gpu.median}\% (\sparkgpu{}); the median between-round CVs
were \stat{cv.between.polars.median}\%, \stat{cv.between.polars_gpu.median}\%,
\stat{cv.between.spark_cpu.median}\% and \stat{cv.between.spark_gpu.median}\%
(Figure~\ref{fig:f5}). Engines whose between-round CV exceeded their
within-session CV: \stat{cv.between_exceeds_within}. Repeating rounds was
still necessary for the intervals, but on this host session-to-session drift
was not larger than trial-to-trial noise.

\paragraph{Order and carryover.} The largest absolute deviation of mean
log-time by session position was \stat{order.position.max_abs_pct}\%, and by
predecessor engine \stat{order.carryover.max_abs_pct}\%
(Figure~\ref{fig:f6}); \stat{order.intervals_including_zero} of the
\stat{order.intervals_total} position and predecessor intervals include zero.
These are descriptive; the Latin square balances but does not remove such
effects, and they are small next to the engine contrasts.

\paragraph{Startup and wall time.} Median launch-to-ready time was
\stat{startup.polars.median_s}\,s, \stat{startup.polars_gpu.median_s}\,s,
\stat{startup.spark_cpu.median_s}\,s and \stat{startup.spark_gpu.median_s}\,s
for the four engines. Startup is outside the query timer but matters for short
jobs: \sparkgpu{}'s median start-up (\stat{startup.spark_gpu.median_s}\,s) was
longer than the sum of its six SF1 baseline query medians
(\stat{legacy.sum_medians.spark_gpu.baseline.sf1}\,s).

\paragraph{Energy proxy.} Mean GPU board power multiplied by query time is
reported in the artifact as a descriptive proxy only; it excludes CPU, memory
and host power and is not an energy measurement.

\paragraph{Earlier A100 study versus H100.}

\begin{figure}[t]
\centering
\studyfig{f09_a100_vs_h100}{0.8\linewidth}
\caption{GPU Spark in the earlier study (\stat{a100_h100.study.label}; saved
evidence of a previous study, not the A100 ladder campaign) versus the H100
campaign, per suite and scale ($t_{\mathrm{A100}}/t_{\mathrm{H100}}$).
Non-equivalent hosts and protocols: \stat{a100_h100.study.difference_headline}
Descriptive only; the A100 ladder campaign (Section~\ref{sec:ladder}) is the
like-for-like comparison. \chartnote}
\label{fig:f9}
\end{figure}

The \sparkgpu{} ratio $t_{\mathrm{A100}}/t_{\mathrm{H100}}$ at SF300 was
\stat{a100_h100.spark_gpu.baseline.sf300.estimate} (baseline) and
\stat{a100_h100.spark_gpu.complex.sf300.estimate} (complex)
(Figure~\ref{fig:f9}). The earlier study (\stat{a100_h100.study.gpu},
\stat{a100_h100.study.vcpus} vCPUs, \stat{a100_h100.study.host_ram_gib}\,GiB
RAM) ran \stat{a100_h100.study.cpu_threads} engine threads against
\stat{a100_h100.study.cpu_threads_h100} here, with
\stat{a100_h100.study.sessions_per_cell} session per cell against
\stat{a100_h100.study.rounds_h100} rounds, and recorded
\stat{a100_h100.study.spark_conf_keys_recorded.a100_study} Spark configuration
keys against \stat{a100_h100.study.spark_conf_keys_recorded.h100_campaign} in
this campaign; the artifact lists every recorded difference. This is
therefore not a GPU-only comparison. The earlier study did not run the Polars
GPU engine, so only \sparkgpu{} is compared, and its CPU engines are not
compared with this campaign.

% ======================================================================
\subsection{The GPU-generation ladder (amendments A2, A3)}
\label{sec:ladder}
% Each host's campaign is analyzed exactly as the H100 campaign; comparisons
% between hosts are descriptive only.
% BEGIN-RAN:any
\ifladderran
Campaigns ran on \stat{hosts.ran_list} (\stat{hosts.n_ran} hosts); hosts not
run: \stat{hosts.not_run_list}. Each campaign is analyzed on its own, exactly
as the H100 campaign above; Appendix~\ref{app:ladder} gives every host's full
contrast table. Table~\ref{tab:ladder} sets the SF300 contrasts side by side,
Table~\ref{tab:ladder-xhost} and Figure~\ref{fig:f11} compare each engine's
times with the H100, and Table~\ref{tab:ladder-coverage} gives GPU coverage per
host. \hostnotrun{} marks a host whose campaign did not run.

\begin{table}[t]
\centering
\caption{Within-host contrasts at SF300 on each host of the ladder:
suite-level geometric-mean speedup ratios with 95\% CIs ($t$, df~=~3,
across that host's four rounds). Each row comes from one host's campaign
alone; rows are not pooled. Crossover: the smallest scale at which the C4
interval lies entirely above~1. \chartnote}
\label{tab:ladder}
\scriptsize
\setlength{\tabcolsep}{3pt}
\begin{tabular}{@{}llccccc@{}}
\toprule
Suite & Host & \shortstack{C1\\\polarscpu/\\\polarsgpu} & \shortstack{C2\\\sparkcpu/\\\sparkgpu} & \shortstack{C3\\\sparkgpu/\\\polarsgpu} & \shortstack{C4\\\polarscpu/\\best GPU} & Crossover \\
\midrule
baseline & A100  & \hostratioci{a100}{c1.baseline.sf300} & \hostratioci{a100}{c2.baseline.sf300} & \hostratioci{a100}{c3.baseline.sf300} & \hostratioci{a100}{c4.baseline.sf300} & \hoststat{a100}{crossover.baseline.c4.scale} \\
         & H100  & \ratioci{c1.baseline.sf300} & \ratioci{c2.baseline.sf300} & \ratioci{c3.baseline.sf300} & \ratioci{c4.baseline.sf300} & \stat{crossover.baseline.c4.scale} \\
         & GH200 & \hostratioci{gh200}{c1.baseline.sf300} & \hostratioci{gh200}{c2.baseline.sf300} & \hostratioci{gh200}{c3.baseline.sf300} & \hostratioci{gh200}{c4.baseline.sf300} & \hoststat{gh200}{crossover.baseline.c4.scale} \\
         & B200  & \hostratioci{b200}{c1.baseline.sf300} & \hostratioci{b200}{c2.baseline.sf300} & \hostratioci{b200}{c3.baseline.sf300} & \hostratioci{b200}{c4.baseline.sf300} & \hoststat{b200}{crossover.baseline.c4.scale} \\
\midrule
complex  & A100  & \hostratioci{a100}{c1.complex.sf300} & \hostratioci{a100}{c2.complex.sf300} & \hostratioci{a100}{c3.complex.sf300} & \hostratioci{a100}{c4.complex.sf300} & \hoststat{a100}{crossover.complex.c4.scale} \\
         & H100  & \ratioci{c1.complex.sf300} & \ratioci{c2.complex.sf300} & \ratioci{c3.complex.sf300} & \ratioci{c4.complex.sf300} & \stat{crossover.complex.c4.scale} \\
         & GH200 & \hostratioci{gh200}{c1.complex.sf300} & \hostratioci{gh200}{c2.complex.sf300} & \hostratioci{gh200}{c3.complex.sf300} & \hostratioci{gh200}{c4.complex.sf300} & \hoststat{gh200}{crossover.complex.c4.scale} \\
         & B200  & \hostratioci{b200}{c1.complex.sf300} & \hostratioci{b200}{c2.complex.sf300} & \hostratioci{b200}{c3.complex.sf300} & \hostratioci{b200}{c4.complex.sf300} & \hoststat{b200}{crossover.complex.c4.scale} \\
\bottomrule
\end{tabular}
\end{table}

\begin{table}[t]
\centering
\caption{Cross-host time ratios at SF300, $t_{h}/t_{\mathrm{H100}}$: suite
geometric mean over queries supported on both hosts, with Welch-type 95\%
intervals on the log scale. Below~1: faster on host~$h$ than on the H100.
Descriptive only: for GPU engines the ratio confounds GPU, CPU, memory,
CPU--GPU link and driver; for CPU engines it reflects a different CPU model
(and, on the GH200, instruction set) at the same 26 threads. \chartnote}
\label{tab:ladder-xhost}
\scriptsize
\setlength{\tabcolsep}{3pt}
\begin{tabular}{@{}llcccc@{}}
\toprule
Suite & Host $h$ & \polarscpu & \polarsgpu & \sparkcpu & \sparkgpu \\
\midrule
baseline & A100  & \xhostratioci{a100}{polars.baseline.sf300} & \xhostratioci{a100}{polars_gpu.baseline.sf300} & \xhostratioci{a100}{spark_cpu.baseline.sf300} & \xhostratioci{a100}{spark_gpu.baseline.sf300} \\
         & GH200 & \xhostratioci{gh200}{polars.baseline.sf300} & \xhostratioci{gh200}{polars_gpu.baseline.sf300} & \xhostratioci{gh200}{spark_cpu.baseline.sf300} & \xhostratioci{gh200}{spark_gpu.baseline.sf300} \\
         & B200  & \xhostratioci{b200}{polars.baseline.sf300} & \xhostratioci{b200}{polars_gpu.baseline.sf300} & \xhostratioci{b200}{spark_cpu.baseline.sf300} & \xhostratioci{b200}{spark_gpu.baseline.sf300} \\
\midrule
complex  & A100  & \xhostratioci{a100}{polars.complex.sf300} & \xhostratioci{a100}{polars_gpu.complex.sf300} & \xhostratioci{a100}{spark_cpu.complex.sf300} & \xhostratioci{a100}{spark_gpu.complex.sf300} \\
         & GH200 & \xhostratioci{gh200}{polars.complex.sf300} & \xhostratioci{gh200}{polars_gpu.complex.sf300} & \xhostratioci{gh200}{spark_cpu.complex.sf300} & \xhostratioci{gh200}{spark_gpu.complex.sf300} \\
         & B200  & \xhostratioci{b200}{polars.complex.sf300} & \xhostratioci{b200}{polars_gpu.complex.sf300} & \xhostratioci{b200}{spark_cpu.complex.sf300} & \xhostratioci{b200}{spark_gpu.complex.sf300} \\
\bottomrule
\end{tabular}
\end{table}

\begin{figure}[t]
\centering
\studyfig{f11_gpu_ladder}{\linewidth}
\caption{The GPU-generation ladder: suite geometric-mean time ratio
$t_{h}/t_{\mathrm{H100}}$ (log axis) for each engine (columns) and suite
(rows), with hosts on the horizontal axis in generation order (A100, H100,
GH200, B200) and the H100 as the reference at~1. Within each host, one marker
per scale factor, SF1 to SF300 from left to right (lighter to darker); bars
are Welch-type 95\% intervals (rounds unpaired). Open markers: fewer queries
compared than at the other hosts (for example after resource exhaustion).
Hosts whose campaign did not run are omitted. Below~1: faster than on the
H100. Descriptive only; the hosts differ in more than the GPU
(Table~\ref{tab:hosts}). \chartnote}
\label{fig:f11}
\end{figure}

\begin{table}[t]
\centering
\caption{GPU coverage per host (amendment A1 classes): query--scale cells
that ran on the GPU with all gates passed, out of all cells, and cells lost
to resource exhaustion (for example out of device memory).}
\label{tab:ladder-coverage}
\scriptsize
\setlength{\tabcolsep}{4pt}
\begin{tabular}{@{}lcccccccc@{}}
\toprule
& \multicolumn{4}{c}{\polarsgpu} & \multicolumn{4}{c}{\sparkgpu} \\
\cmidrule(lr){2-5}\cmidrule(l){6-9}
& \multicolumn{2}{c}{ran / total} & \multicolumn{2}{c}{exhausted} & \multicolumn{2}{c}{ran / total} & \multicolumn{2}{c}{exhausted} \\
Host & baseline & complex & baseline & complex & baseline & complex & baseline & complex \\
\midrule
A100  & \hoststat{a100}{coverage.polars_gpu.baseline.supported} / \hoststat{a100}{coverage.polars_gpu.baseline.total} & \hoststat{a100}{coverage.polars_gpu.complex.supported} / \hoststat{a100}{coverage.polars_gpu.complex.total} & \hoststat{a100}{coverage.polars_gpu.baseline.resource_exhaustion} & \hoststat{a100}{coverage.polars_gpu.complex.resource_exhaustion} & \hoststat{a100}{coverage.spark_gpu.baseline.supported} / \hoststat{a100}{coverage.spark_gpu.baseline.total} & \hoststat{a100}{coverage.spark_gpu.complex.supported} / \hoststat{a100}{coverage.spark_gpu.complex.total} & \hoststat{a100}{coverage.spark_gpu.baseline.resource_exhaustion} & \hoststat{a100}{coverage.spark_gpu.complex.resource_exhaustion} \\
H100  & \stat{coverage.polars_gpu.baseline.supported} / \stat{coverage.polars_gpu.baseline.total} & \stat{coverage.polars_gpu.complex.supported} / \stat{coverage.polars_gpu.complex.total} & \stat{coverage.polars_gpu.baseline.resource_exhaustion} & \stat{coverage.polars_gpu.complex.resource_exhaustion} & \stat{coverage.spark_gpu.baseline.supported} / \stat{coverage.spark_gpu.baseline.total} & \stat{coverage.spark_gpu.complex.supported} / \stat{coverage.spark_gpu.complex.total} & \stat{coverage.spark_gpu.baseline.resource_exhaustion} & \stat{coverage.spark_gpu.complex.resource_exhaustion} \\
GH200 & \hoststat{gh200}{coverage.polars_gpu.baseline.supported} / \hoststat{gh200}{coverage.polars_gpu.baseline.total} & \hoststat{gh200}{coverage.polars_gpu.complex.supported} / \hoststat{gh200}{coverage.polars_gpu.complex.total} & \hoststat{gh200}{coverage.polars_gpu.baseline.resource_exhaustion} & \hoststat{gh200}{coverage.polars_gpu.complex.resource_exhaustion} & \hoststat{gh200}{coverage.spark_gpu.baseline.supported} / \hoststat{gh200}{coverage.spark_gpu.baseline.total} & \hoststat{gh200}{coverage.spark_gpu.complex.supported} / \hoststat{gh200}{coverage.spark_gpu.complex.total} & \hoststat{gh200}{coverage.spark_gpu.baseline.resource_exhaustion} & \hoststat{gh200}{coverage.spark_gpu.complex.resource_exhaustion} \\
B200  & \hoststat{b200}{coverage.polars_gpu.baseline.supported} / \hoststat{b200}{coverage.polars_gpu.baseline.total} & \hoststat{b200}{coverage.polars_gpu.complex.supported} / \hoststat{b200}{coverage.polars_gpu.complex.total} & \hoststat{b200}{coverage.polars_gpu.baseline.resource_exhaustion} & \hoststat{b200}{coverage.polars_gpu.complex.resource_exhaustion} & \hoststat{b200}{coverage.spark_gpu.baseline.supported} / \hoststat{b200}{coverage.spark_gpu.baseline.total} & \hoststat{b200}{coverage.spark_gpu.complex.supported} / \hoststat{b200}{coverage.spark_gpu.complex.total} & \hoststat{b200}{coverage.spark_gpu.baseline.resource_exhaustion} & \hoststat{b200}{coverage.spark_gpu.complex.resource_exhaustion} \\
\bottomrule
\end{tabular}
\end{table}

The H100's qualitative pattern recurred on every host that ran
(Table~\ref{tab:ladder}). C4's baseline interval lay below~1
\stat{ladder.c4.baseline.below_text}, and above~1 from each host's crossover
scale onwards; on the complex suite it lay above~1
\stat{ladder.c4.complex.above_text}. \polarsgpu{} was faster
than \sparkgpu{} in \stat{ladder.c3.cells_above} of \stat{ladder.c3.cells}
suite--scale cells across the ladder.

Across hosts (Table~\ref{tab:ladder-xhost}, Figure~\ref{fig:f11}):
% BEGIN-RAN:a100
\ifaonehundredran
the GPU engines were slower on the A100 than on the H100, with
$t_{\mathrm{A100}}/t_{\mathrm{H100}}$ from \stat{ladder.xhost_gpu.a100.min} to
\stat{ladder.xhost_gpu.a100.max};
\fi
% END-RAN:a100
on the \stat{ladder.xhost_gpu.later.hosts} the GPU engines' ratios ranged from
\stat{ladder.xhost_gpu.later.min} to \stat{ladder.xhost_gpu.later.max}. The CPU
engines followed a different pattern from the GPU engines:
% BEGIN-RAN:gh200
\ifghtwohundredran
\polarscpu{} was faster on the GH200 (ratios from
\stat{xhost.gh200.polars.overall.min} to \stat{xhost.gh200.polars.overall.max}),
\fi
% END-RAN:gh200
% BEGIN-RAN:a100
\ifaonehundredran
slower on the A100 (\stat{xhost.a100.polars.overall.min} to
\stat{xhost.a100.polars.overall.max}),
\fi
% END-RAN:a100
which is consistent with the hosts' different CPUs and with the number of
physical cores the 26 threads occupy (\hoststat{gh200}{host.cpus_arg_cores_text}
on the GH200 against \stat{host.cpus_arg_cores_text} on the H100;
Table~\ref{tab:hosts}). We report these differences as observed and do not
attribute them to a cause; each rung changes the whole host
(Section~\ref{sec:threats}).

\begin{table}[t]
\centering
\caption{GPU telemetry at SF300 (\stat{secondary.gpu_telemetry.meta.label}).
Pooled nvidia-smi samples every \stat{secondary.gpu_telemetry.meta.sample_interval_ms}\,ms
inside measured actions of at least \stat{secondary.gpu_telemetry.meta.min_action_seconds}\,s.
Util: median \texttt{utilization.gpu} (share of the sample period in which a
kernel was executing, not SM occupancy), \%. Idle: share of samples below
\stat{secondary.gpu_telemetry.meta.idle_threshold_pct}\% utilization, \%. Power:
median board power as a share of the card's power limit, \%. Mem: median
memory-controller utilization (\texttt{utilization.memory}), \%. The median SM
clock was at the card's maximum in \stat{ladder.gpu_tel.clock_at_max} of
\stat{ladder.gpu_tel.cells} populated suite--scale--engine cells.}
\label{tab:ladder-telemetry}
\scriptsize
\setlength{\tabcolsep}{4pt}
\begin{tabular}{@{}llrrrrrrrr@{}}
\toprule
& & \multicolumn{4}{c}{baseline} & \multicolumn{4}{c}{complex} \\
\cmidrule(lr){3-6}\cmidrule(l){7-10}
Host & Engine & Util & Idle & Power & Mem & Util & Idle & Power & Mem \\
\midrule
A100  & \polarsgpu & \hoststat{a100}{gpu_tel.polars_gpu.baseline.sf300.util} & \hoststat{a100}{gpu_tel.polars_gpu.baseline.sf300.idle} & \hoststat{a100}{gpu_tel.polars_gpu.baseline.sf300.power} & \hoststat{a100}{gpu_tel.polars_gpu.baseline.sf300.memutil} & \hoststat{a100}{gpu_tel.polars_gpu.complex.sf300.util} & \hoststat{a100}{gpu_tel.polars_gpu.complex.sf300.idle} & \hoststat{a100}{gpu_tel.polars_gpu.complex.sf300.power} & \hoststat{a100}{gpu_tel.polars_gpu.complex.sf300.memutil} \\
      & \sparkgpu & \hoststat{a100}{gpu_tel.spark_gpu.baseline.sf300.util} & \hoststat{a100}{gpu_tel.spark_gpu.baseline.sf300.idle} & \hoststat{a100}{gpu_tel.spark_gpu.baseline.sf300.power} & \hoststat{a100}{gpu_tel.spark_gpu.baseline.sf300.memutil} & \hoststat{a100}{gpu_tel.spark_gpu.complex.sf300.util} & \hoststat{a100}{gpu_tel.spark_gpu.complex.sf300.idle} & \hoststat{a100}{gpu_tel.spark_gpu.complex.sf300.power} & \hoststat{a100}{gpu_tel.spark_gpu.complex.sf300.memutil} \\
\addlinespace
H100  & \polarsgpu & \stat{gpu_tel.polars_gpu.baseline.sf300.util} & \stat{gpu_tel.polars_gpu.baseline.sf300.idle} & \stat{gpu_tel.polars_gpu.baseline.sf300.power} & \stat{gpu_tel.polars_gpu.baseline.sf300.memutil} & \stat{gpu_tel.polars_gpu.complex.sf300.util} & \stat{gpu_tel.polars_gpu.complex.sf300.idle} & \stat{gpu_tel.polars_gpu.complex.sf300.power} & \stat{gpu_tel.polars_gpu.complex.sf300.memutil} \\
      & \sparkgpu & \stat{gpu_tel.spark_gpu.baseline.sf300.util} & \stat{gpu_tel.spark_gpu.baseline.sf300.idle} & \stat{gpu_tel.spark_gpu.baseline.sf300.power} & \stat{gpu_tel.spark_gpu.baseline.sf300.memutil} & \stat{gpu_tel.spark_gpu.complex.sf300.util} & \stat{gpu_tel.spark_gpu.complex.sf300.idle} & \stat{gpu_tel.spark_gpu.complex.sf300.power} & \stat{gpu_tel.spark_gpu.complex.sf300.memutil} \\
\addlinespace
GH200 & \polarsgpu & \hoststat{gh200}{gpu_tel.polars_gpu.baseline.sf300.util} & \hoststat{gh200}{gpu_tel.polars_gpu.baseline.sf300.idle} & \hoststat{gh200}{gpu_tel.polars_gpu.baseline.sf300.power} & \hoststat{gh200}{gpu_tel.polars_gpu.baseline.sf300.memutil} & \hoststat{gh200}{gpu_tel.polars_gpu.complex.sf300.util} & \hoststat{gh200}{gpu_tel.polars_gpu.complex.sf300.idle} & \hoststat{gh200}{gpu_tel.polars_gpu.complex.sf300.power} & \hoststat{gh200}{gpu_tel.polars_gpu.complex.sf300.memutil} \\
      & \sparkgpu & \hoststat{gh200}{gpu_tel.spark_gpu.baseline.sf300.util} & \hoststat{gh200}{gpu_tel.spark_gpu.baseline.sf300.idle} & \hoststat{gh200}{gpu_tel.spark_gpu.baseline.sf300.power} & \hoststat{gh200}{gpu_tel.spark_gpu.baseline.sf300.memutil} & \hoststat{gh200}{gpu_tel.spark_gpu.complex.sf300.util} & \hoststat{gh200}{gpu_tel.spark_gpu.complex.sf300.idle} & \hoststat{gh200}{gpu_tel.spark_gpu.complex.sf300.power} & \hoststat{gh200}{gpu_tel.spark_gpu.complex.sf300.memutil} \\
\addlinespace
B200  & \polarsgpu & \hoststat{b200}{gpu_tel.polars_gpu.baseline.sf300.util} & \hoststat{b200}{gpu_tel.polars_gpu.baseline.sf300.idle} & \hoststat{b200}{gpu_tel.polars_gpu.baseline.sf300.power} & \hoststat{b200}{gpu_tel.polars_gpu.baseline.sf300.memutil} & \hoststat{b200}{gpu_tel.polars_gpu.complex.sf300.util} & \hoststat{b200}{gpu_tel.polars_gpu.complex.sf300.idle} & \hoststat{b200}{gpu_tel.polars_gpu.complex.sf300.power} & \hoststat{b200}{gpu_tel.polars_gpu.complex.sf300.memutil} \\
      & \sparkgpu & \hoststat{b200}{gpu_tel.spark_gpu.baseline.sf300.util} & \hoststat{b200}{gpu_tel.spark_gpu.baseline.sf300.idle} & \hoststat{b200}{gpu_tel.spark_gpu.baseline.sf300.power} & \hoststat{b200}{gpu_tel.spark_gpu.baseline.sf300.memutil} & \hoststat{b200}{gpu_tel.spark_gpu.complex.sf300.util} & \hoststat{b200}{gpu_tel.spark_gpu.complex.sf300.idle} & \hoststat{b200}{gpu_tel.spark_gpu.complex.sf300.power} & \hoststat{b200}{gpu_tel.spark_gpu.complex.sf300.memutil} \\
\bottomrule
\end{tabular}
\end{table}

\begin{figure}[t]
\centering
\studyfig{f12_gpu_idle_fraction}{0.8\linewidth}
\caption{Share of GPU-idle samples (utilization below
\stat{secondary.gpu_telemetry.meta.idle_threshold_pct}\%) inside measured
actions of at least \stat{secondary.gpu_telemetry.meta.min_action_seconds}\,s,
by scale, GPU engine and host (\stat{secondary.gpu_telemetry.meta.label}).
Cells with no qualifying action are omitted.}
\label{fig:f12}
\end{figure}

GPU telemetry (Table~\ref{tab:ladder-telemetry}, Figure~\ref{fig:f12};
\stat{secondary.gpu_telemetry.meta.label}) describes how busy the GPU was during
measured actions; it is not part of any estimate. At SF300 on the H100, the
median utilization was \stat{gpu_tel.polars_gpu.complex.sf300.util}\% for
\polarsgpu{} and \stat{gpu_tel.spark_gpu.complex.sf300.util}\% for \sparkgpu{} on
the complex suite, with GPU-idle shares of
\stat{gpu_tel.polars_gpu.complex.sf300.idle}\% and
\stat{gpu_tel.spark_gpu.complex.sf300.idle}\%. The median SM clock was at the
card's maximum in \stat{ladder.gpu_tel.clock_at_max} of \stat{ladder.gpu_tel.cells}
populated cells. Utilization here is the share of time a kernel was running, not
how much of the GPU it used, and we draw no causal conclusion from it.

% BEGIN-RAN:a100
\ifaonehundredran
\paragraph{A100.} The A100 campaign completed
\stat{a100.meta.rounds_completed} rounds (\stat{a100.meta.rounds_excluded_count} excluded,
\stat{a100.meta.rounds_repeated} repeated; input gate \stat{a100.meta.input_gate}; design
version \texttt{\stat{a100.meta.design_commit_short}} with
\stat{a100.meta.amendments} amendments). On the
baseline suite at SF300, the time ratio $t_{\mathrm{A100}}/t_{\mathrm{H100}}$
was \ratiotext{xhost.a100.polars_gpu.baseline.sf300} for \polarsgpu{} and
\ratiotext{xhost.a100.spark_gpu.baseline.sf300} for \sparkgpu{}. Peak sampled
device memory on the complex suite at SF300 reached
\stat{a100.peak_gpu_mem_gib.polars_gpu.complex.sf300.max}\,GiB for \polarsgpu{}
(sampled device memory; on the A100, \stat{a100.gpu_mem_note.polars_gpu}) and
\stat{a100.task_gpu_mem_gib.spark_gpu.complex.sf300.max}\,GiB for \sparkgpu{}
(task-reported maximum), against
\stat{a100.host.gpu_memory_gib}\,GiB of device memory.
Resource exhaustion removed \stat{a100.coverage.polars_gpu.baseline.resource_exhaustion}
baseline and \stat{a100.coverage.polars_gpu.complex.resource_exhaustion} complex
\polarsgpu{} cells and \stat{a100.coverage.spark_gpu.complex.resource_exhaustion}
complex \sparkgpu{} cells on the A100. With the Polars GPU engine included,
the A100's C4 at SF300 was \ratiotext{a100.c4.baseline.sf300} (baseline) and
\ratiotext{a100.c4.complex.sf300} (complex), with the crossover at
\stat{a100.crossover.c4.text}; the earlier three-engine study, which
had no Polars GPU engine, could not measure this contrast.
\fi
% END-RAN:a100
% BEGIN-RAN:gh200
\ifghtwohundredran
\paragraph{GH200.} The GH200 campaign completed
\stat{gh200.meta.rounds_completed} rounds (\stat{gh200.meta.rounds_excluded_count} excluded,
\stat{gh200.meta.rounds_repeated} repeated; input gate \stat{gh200.meta.input_gate}; design
version \texttt{\stat{gh200.meta.design_commit_short}} with
\stat{gh200.meta.amendments} amendments) on aarch64.
On the baseline suite at SF300, the time ratio
$t_{\mathrm{GH200}}/t_{\mathrm{H100}}$ of the CPU engines, running on the
Grace CPU against the H100 host's x86\_64 CPU at the same 26 threads, was
\ratiotext{xhost.gh200.polars.baseline.sf300} for \polarscpu{} and
\ratiotext{xhost.gh200.spark_cpu.baseline.sf300} for \sparkcpu{}; for the
GPU engines it was \ratiotext{xhost.gh200.polars_gpu.baseline.sf300}
(\polarsgpu{}) and \ratiotext{xhost.gh200.spark_gpu.baseline.sf300}
(\sparkgpu{}). The CPU and GPU sides moved differently: \polarscpu{}'s
interval lay below~1 (faster on the GH200)
\stat{xhost.gh200.polars.both.below_text}, while
\polarsgpu{} ranged from \stat{xhost.gh200.polars_gpu.overall.min} to
\stat{xhost.gh200.polars_gpu.overall.max}. Because C1 and C4 use each host's
own CPU engine, the GH200's C4 at SF300 was
\stat{gh200.c4.baseline.sf300.vs_h100} on the baseline suite and
\stat{gh200.c4.complex.sf300.vs_h100} on the complex suite.
\fi
% END-RAN:gh200
% BEGIN-RAN:b200
\ifbtwohundredran
\paragraph{B200.} The B200 campaign completed
\stat{b200.meta.rounds_completed} rounds (\stat{b200.meta.rounds_excluded_count} excluded,
\stat{b200.meta.rounds_repeated} repeated; input gate \stat{b200.meta.input_gate}; design
version \texttt{\stat{b200.meta.design_commit_short}} with
\stat{b200.meta.amendments} amendments). At SF300
the time ratio $t_{\mathrm{B200}}/t_{\mathrm{H100}}$ was
\ratiotext{xhost.b200.polars_gpu.baseline.sf300} for \polarsgpu{} and
\ratiotext{xhost.b200.spark_gpu.baseline.sf300} for \sparkgpu{} on the
baseline suite, and \ratiotext{xhost.b200.polars_gpu.complex.sf300} and
\ratiotext{xhost.b200.spark_gpu.complex.sf300} on the complex suite.
For the \sparkgpu{} device-memory sample on the B200,
\stat{b200.gpu_mem_note.spark_gpu}; on the H100,
\stat{gpu_mem_note.spark_gpu}. Where the sample is not constant in every
session, Figure~\ref{fig:f7}'s convention omits the pool-reservation line. We
report this as observed; the task-reported maxima do not depend on it.
\polarsgpu{} on the complex suite was slower on the B200 than on the H100
(interval above~1 at \stat{xhost.b200.polars_gpu.complex.summary.above}, up to
\stat{xhost.b200.polars_gpu.complex.summary.max}), while \sparkgpu{} ranged from
\stat{xhost.b200.spark_gpu.overall.min} to \stat{xhost.b200.spark_gpu.overall.max}.
The B200's C4 at SF300 was \stat{b200.c4.baseline.sf300.vs_h100} on the
baseline suite and \stat{b200.c4.complex.sf300.vs_h100} on the complex suite,
with the crossover at \stat{b200.crossover.c4.text}. These timings alone do not explain the
differences between the B200 and the H100.
\fi
% END-RAN:b200
\fi
% END-RAN:any
% BEGIN-NOTRUN:any
\ifladderran\else
None of the additional ladder campaigns ran (reasons in
Section~\ref{sec:ladder-design}). As pre-registered, no software versions
were changed to make them run, so all results in this paper come from the H100
campaign.
\fi
% END-NOTRUN:any

% ======================================================================
\subsection{Exploratory analyses (not pre-registered)}
\label{sec:exploratory}
The analyses in this subsection were \emph{not} pre-registered. They were
added after the side study showed how much the Polars CPU version matters
(Section~\ref{sec:rq5}), and they are reported only here: they do not enter
Table~\ref{tab:contrasts}, the ladder tables, the RQ3 crossover rule or any
primary claim. E1 compares Polars~2.0.0 on the CPU with \polarsgpu{} (which
runs on Polars~1.44.2), and E4 compares Polars~2.0.0 on the CPU with the best
GPU engine per query. Because the two Polars versions come from different
sessions (side study versus main rounds), these ratios also combine two
measurement protocols.

\begin{table}[t]
\centering
\caption{\textsc{Exploratory}, \stat{exploratory.e1.meta.label}: E1 and E4
next to the pre-registered C1 and C4. E1 = $t$(Polars \stat{exploratory.e1.meta.cpu_version}
CPU, side-study sessions) / $t$(\polarsgpu{}, main rounds); E4 uses the
faster GPU engine per query and round instead. Above~1: the GPU side was
faster. C1 and C4: paired $t$-interval across rounds (df~=~3). E1 and E4:
\stat{exploratory_summary.interval_note}. Not part of any primary result or
the crossover rule. \chartnote}
\label{tab:exploratory}
\scriptsize
\setlength{\tabcolsep}{3pt}
\begin{tabular}{@{}llcccc@{}}
\toprule
Suite & Scale & \shortstack{C1\\(pre-registered)} & \shortstack{E1\\(exploratory)} & \shortstack{C4\\(pre-registered)} & \shortstack{E4\\(exploratory)} \\
\midrule
baseline & SF1   & \ratioci{c1.baseline.sf1} & \ratioci{exploratory.e1.baseline.sf1} & \ratioci{c4.baseline.sf1} & \ratioci{exploratory.e4.baseline.sf1} \\
         & SF10  & \ratioci{c1.baseline.sf10} & \ratioci{exploratory.e1.baseline.sf10} & \ratioci{c4.baseline.sf10} & \ratioci{exploratory.e4.baseline.sf10} \\
         & SF30  & \ratioci{c1.baseline.sf30} & \ratioci{exploratory.e1.baseline.sf30} & \ratioci{c4.baseline.sf30} & \ratioci{exploratory.e4.baseline.sf30} \\
         & SF100 & \ratioci{c1.baseline.sf100} & \ratioci{exploratory.e1.baseline.sf100} & \ratioci{c4.baseline.sf100} & \ratioci{exploratory.e4.baseline.sf100} \\
         & SF300 & \ratioci{c1.baseline.sf300} & \ratioci{exploratory.e1.baseline.sf300} & \ratioci{c4.baseline.sf300} & \ratioci{exploratory.e4.baseline.sf300} \\
\midrule
complex  & SF30  & \ratioci{c1.complex.sf30} & \ratioci{exploratory.e1.complex.sf30} & \ratioci{c4.complex.sf30} & \ratioci{exploratory.e4.complex.sf30} \\
         & SF100 & \ratioci{c1.complex.sf100} & \ratioci{exploratory.e1.complex.sf100} & \ratioci{c4.complex.sf100} & \ratioci{exploratory.e4.complex.sf100} \\
         & SF300 & \ratioci{c1.complex.sf300} & \ratioci{exploratory.e1.complex.sf300} & \ratioci{c4.complex.sf300} & \ratioci{exploratory.e4.complex.sf300} \\
\bottomrule
\end{tabular}
\end{table}

\begin{figure}[t]
\centering
\studyfig{f02x_c1_e1_exploratory}{\linewidth}
\caption{\textsc{Exploratory} (post-hoc; not in the locked design): suite
geometric-mean ratios versus scale for C1 and C4 (pre-registered) and E1 and
E4 (exploratory, against Polars~2.0.0 on the CPU). The exploratory panels are
marked as such in the figure; their bands are Welch-type 95\% intervals.
\chartnote}
\label{fig:f2x}
\end{figure}

At SF300, E1 (exploratory, post-hoc) was
\ratiotext{exploratory.e1.baseline.sf300} on the baseline suite and
\ratiotext{exploratory.e1.complex.sf300} on the complex suite, against C1 of
\ratiotext{c1.baseline.sf300} and \ratiotext{c1.complex.sf300}; E4
(exploratory, post-hoc) was \ratiotext{exploratory.e4.baseline.sf300} and
\ratiotext{exploratory.e4.complex.sf300}, against C4 of
\ratiotext{c4.baseline.sf300} and \ratiotext{c4.complex.sf300}
(Table~\ref{tab:exploratory}, Figure~\ref{fig:f2x}). E1 and E4 use
\stat{exploratory.e1.complex.sf300.n_queries} complex queries at SF300 for E1
and \stat{exploratory.e4.complex.sf300.n_queries} for E4 (queries the GPU side
could not run are excluded, as in C1 and C4). E1 was smaller than C1
\stat{exploratory.e1.both.below_c1_text}: against
Polars~2.0.0 on the CPU, \polarsgpu{}'s advantage shrank. E1's interval still
lay above~1 \stat{exploratory.e1.both.above_text}, and below~1 at
\stat{exploratory.e1.baseline.summary.below} on the baseline suite. These are
exploratory (post-hoc) results; no crossover statement is drawn from E1 or
E4.
% NOTE: every E number must keep its "exploratory" label.

% ======================================================================
\section{Discussion}
\label{sec:discussion}

\paragraph{When does a GPU help?} On this workload the answer depends on scale
and query shape. At the smallest scale both GPU engines lost to their own
framework's CPU engine (C1 below~1 at \stat{c1.baseline.summary.below}, C2 at
\stat{c2.baseline.summary.below}), which is consistent with fixed per-query
costs that a small input cannot amortize. \polarsgpu{} was faster than Polars
1.44.2 on the CPU \stat{c1.both.above_text}. The largest per-query gains were on grouped aggregations
and joins and the smallest on the selective scan-and-sum \q{q06}, whose
SF300 interval \stat{c1.baseline.sf300.q.q06.verdict}. That pattern matches the
Polars guide's statement that grouped aggregations and joins benefit most and
I/O-bound queries least~\citep{polars_gpu_guide}, although our design does not
isolate query shape as a factor.

\paragraph{GPU versus a strong CPU baseline.} C4 asks whether the faster GPU
engine beats Polars on the CPU of the same host, a different question from the
within-framework speedups C1 and C2. On the H100 it did so from
\stat{crossover.baseline.c4.scale} (baseline) and
\stat{crossover.complex.c4.scale} (complex), reaching
\ratiotext{c4.baseline.sf300} and \ratiotext{c4.complex.sf300} at SF300. These
are multiples of a few, not orders of magnitude, in line with the expectation
of \citet{shanbhag2020gpucpu} that well-implemented CPU code bounds GPU speedups
near the bandwidth ratio; with Parquet reads and device-to-host transfers inside
the timer, the transfer costs analyzed by \citet{lutz2020pump} are part of
every GPU time here.

\paragraph{A version-locked CPU baseline.} The Polars GPU engine is tied to
Polars 1.44.x, and Polars 2.0.0 on the CPU was \stat{bridge_summary.both}. The
exploratory E1 (Section~\ref{sec:exploratory}) shows the GPU's advantage
shrinking against the newer release but, as an exploratory and post-hoc result,
still above~1 from \stat{exploratory.e1.baseline.summary.above} on the baseline
suite. The largest single C4 value at SF300,
\ratiotext{c4.complex.sf300.q.cx_window} for \q{cx\_window} (run by \sparkgpu{},
since \polarsgpu{} rejects it as unsupported), is on the query where Polars 2.0.0 gained
most on the CPU (\q{\stat{bridge_summary.complex.extreme_query}},
\stat{bridge_summary.complex.extreme.min}--\stat{bridge_summary.complex.extreme.max}).
This value comes largely from the slow Polars~1.44.2 CPU path and is not a
measure of GPU performance. The version lock is itself a practical cost of
choosing the GPU engine: it ties the CPU side to an older Polars.

\paragraph{Two routes to the GPU.} \polarsgpu{} was faster than \sparkgpu{} in
\stat{ladder.c3.cells_above} of \stat{ladder.c3.cells} suite--scale cells across
the hosts (C3 from \stat{ladder.c3.min} to \stat{ladder.c3.max}). Both engines
execute through libcudf (in different releases), so the gap may reflect what
surrounds the device work: Spark's JVM, task scheduling, shuffle and plan
translation, which are inside our timer, against a single-process Polars plan.
We did not measure these components separately. The descriptive telemetry is
compatible with, but does not establish, such an overhead: on the H100's
complex suite at SF300, \sparkgpu{}'s samples were GPU-idle
\stat{gpu_tel.spark_gpu.complex.sf300.idle}\% of the time against
\stat{gpu_tel.polars_gpu.complex.sf300.idle}\% for \polarsgpu{}
(\stat{secondary.gpu_telemetry.meta.label}). \sparkgpu{} has the
coverage advantage: it ran every complex query, including the two that
\polarsgpu{} rejects (RQ4), and won most of the complex cells \polarsgpu{} did
not win (\stat{fastest.spark_gpu.complex.won_text}).

\paragraph{Practical guidance.} Within the limits of this study: for data at or
above the crossover scale (here \stat{crossover.baseline.c4.scale} on the
H100's baseline suite) on a host like these, try the Polars GPU engine first
and run it once with \texttt{raise\_on\_fail=True} to find queries it would
silently hand back to the CPU; for small inputs Polars on the CPU was faster;
compare against the newest Polars on the CPU, not only the version the GPU
engine supports; and if a query is not supported by the Polars GPU engine,
Spark with the cuDF plugin may run it on the GPU.

% BEGIN-RAN:any
\ifladderran
\paragraph{Does a newer GPU help?} Two questions need separating. First, does
the same engine get faster on a newer GPU (Table~\ref{tab:ladder-xhost},
Figure~\ref{fig:f11})? From the A100 to the H100, yes, for both GPU engines
($t_{\mathrm{A100}}/t_{\mathrm{H100}}$ from \stat{ladder.xhost_gpu.a100.min} to
\stat{ladder.xhost_gpu.a100.max}). Beyond the H100, the GPU engines' ratios on
the GH200 and B200 ranged from \stat{ladder.xhost_gpu.later.min} to
\stat{ladder.xhost_gpu.later.max}, and \polarsgpu{} on the complex suite was
slower on the B200. Second, does a newer GPU change the decision, that is, C4
on that host (Table~\ref{tab:ladder})? Not qualitatively: every host showed the
same pattern of CPU wins at the smallest scale and GPU wins at the largest. A
faster GPU engine need not raise C4 if the host's CPU engine also changes, as
on the GH200. Why the GPU engines did not speed up beyond the H100 is not
answered by these timings. The GPU telemetry
(Table~\ref{tab:ladder-telemetry}, Figure~\ref{fig:f12}) is descriptive and not
pre-registered (\stat{secondary.gpu_telemetry.meta.label}). It pools nvidia-smi
samples taken every \stat{secondary.gpu_telemetry.meta.sample_interval_ms}\,ms and
leaves out actions shorter than \stat{secondary.gpu_telemetry.meta.min_action_seconds}\,s,
so \polarsgpu{} has qualifying cells only from
\stat{gpu_tel.ladder.polars_gpu.smallest_scale} upward and none at SF1. Its
utilization measure, \texttt{utilization.gpu}, is the share of time any kernel
was running, not SM occupancy or throughput. At SF300, \polarsgpu{}'s GPU-idle
share was at most \stat{gpu_tel.ladder.polars_gpu.sf300.idle.max}\% on every
host, while its median power was only \stat{gpu_tel.ladder.polars_gpu.sf300.power.min}\%
to \stat{gpu_tel.ladder.polars_gpu.sf300.power.max}\% of the board limit and its
median memory-controller utilization \stat{gpu_tel.ladder.polars_gpu.sf300.memutil.min}\%
to \stat{gpu_tel.ladder.polars_gpu.sf300.memutil.max}\%, with SM clocks at their
maximum. This is consistent with \polarsgpu{} work that keeps the device
occupied without saturating power or memory bandwidth, which would limit the
gains from a faster GPU. \sparkgpu{}'s GPU-idle share at SF300 ranged from
\stat{gpu_tel.ladder.spark_gpu.sf300.idle.min}\% to
\stat{gpu_tel.ladder.spark_gpu.sf300.idle.max}\% (highest:
\stat{gpu_tel.ladder.spark_gpu.sf300.idle.max_where}), which is consistent with
\sparkgpu{} waiting on host-side work. Each rung changes the whole host, so we
do not attribute the timings to any one component.

\paragraph{Is the H100 picture robust?} For RQ1 to RQ4 the direction of the
H100 conclusions recurred on the other hosts: C4 below~1 at SF1 and above~1 at
the larger scales, \polarsgpu{} ahead of \sparkgpu{} everywhere, and coverage
limited to the pre-declared gaps. Individual intervals near the crossover
differ (for example, some C2 intervals at the smallest scales include~1 on
other hosts; Appendix~\ref{app:ladder}). The
magnitudes differ between hosts (Table~\ref{tab:ladder}), and some crossovers
moved by one scale step. Because every campaign used the same author, code and
software, this is robustness to hardware, not independent replication.
% BEGIN-RAN:gh200
\ifghtwohundredran

\paragraph{The Grace~Hopper rung.} The GH200 pairs a Hopper GPU with an Arm CPU
and a coherent CPU--GPU link. \polarscpu{} was faster there than on the H100
(ratios from \stat{xhost.gh200.polars.overall.min} to
\stat{xhost.gh200.polars.overall.max}), with 26 threads on
\stat{gh200.host.cpus_arg_cores_text} physical cores against
\stat{host.cpus_arg_cores_text} on the H100, while its GPU engines stayed within
\stat{ladder.xhost_gpu.gh200.min}--\stat{ladder.xhost_gpu.gh200.max} of the
H100's. The faster CPU baseline lowers C1 and C4 on this host. A faster
CPU--GPU link would be expected to help transfer-bound queries
\citep{lutz2020pump}; the GPU engines' small ratios do not show a large effect,
but our design cannot separate the link from the other host differences.
\fi
% END-RAN:gh200
% BEGIN-RAN:a100
\ifaonehundredran

\paragraph{The four-engine A100 result.} The earlier A100 study found that
Polars CPU kept the lower aggregate time at every observed scale, with GPU
Spark winning only selected queries. With the Polars GPU engine added, the
A100 campaign shows the best GPU engine ahead of Polars CPU from
\stat{a100.crossover.c4.text}, so the earlier conclusion was
specific to the engines it compared. The 40\,GB device did not cost coverage
(Table~\ref{tab:ladder-coverage}), although \polarsgpu{}'s memory pool came
close to its capacity at complex SF300 (Section~\ref{sec:ladder}).
\fi
% END-RAN:a100
\fi
% END-RAN:any

\paragraph{Comparison with the earlier A100 study.} GPU Spark was faster on the
H100 than in the earlier study ($t_{\mathrm{A100}}/t_{\mathrm{H100}}$
\stat{a100_h100.spark_gpu.baseline.sf300.estimate} and
\stat{a100_h100.spark_gpu.complex.sf300.estimate} at SF300), but that
comparison mixes hosts, thread counts and protocols (Figure~\ref{fig:f9}). The
A100 ladder campaign is the like-for-like A100 result, and it is the one we use
for A100 statements.

% ======================================================================
\section{Threats to Validity}
\label{sec:threats}

\subsection{Internal validity}
\begin{itemize}[leftmargin=*,itemsep=2pt]
  \item \textbf{Few rounds.} Four rounds give $\text{df}=3$ and a $t$
    multiplier of about 3.18, so intervals are wide. A non-overlapping
    interval is informative; an interval including~1 does not show
    equivalence.
  \item \textbf{Order and carryover.} The Williams square balances position
    and first-order carryover across engines but does not remove them;
    higher-order carryover is not balanced. We report position and predecessor
    effects descriptively (Figure~\ref{fig:f6}).
  \item \textbf{Within-session trials are not independent.} Three trials share
    one process, JIT and memory-pool state; we use their median as one
    observation per round rather than as replicates.
  \item \textbf{Unlocked clocks.} GPU and CPU clocks were not locked; they and
    the temperatures were recorded at 200\,ms. Thermal or power throttling would
    show up as between-round variation; it is not controlled.
  \item \textbf{Page cache.} Files are preheated before every session, but
    residency is not verified; at SF300 the inputs approach half of host RAM
    and compete with engine memory.
  \item \textbf{Garbage-collection barriers} in the complex suite change the
    heap state each query sees; they apply equally to all engines but differ
    in cost (Polars versus JVM).
  \item \textbf{Adaptive plans vary between trials.} Spark's adaptive query
    execution can end the same query in different final plans in different
    trials. On the H100, \stat{plan_variability.cells_variable} of
    \stat{plan_variability.cells_checked} Spark query--scale cells had trials
    ending in more than one final plan (\stat{plan_variability.trials_with_plan}
    trials with a saved plan; the cells are listed in the artifact). This is a
    statement about trials, not a defect: every measured trial enters its
    cell's median as designed, so in such a cell the median also reflects which
    plan adaptive execution chose, and its spread is larger. On the H100 the
    variable cell was \stat{plan_variability.cells_text}; the other hosts' are
    listed in the artifact.
% BEGIN-RAN:any
\ifladderran
    The same Spark GPU cell varied on the other hosts as well
    (\hoststat{a100}{plan_variability.cells_text} on the A100;
    \hoststat{gh200}{plan_variability.cells_text} on the GH200;
    \hoststat{b200}{plan_variability.cells_text} on the B200).
\fi
% END-RAN:any
  \item \textbf{Selecting the best GPU engine for C4.} In each round, C4
    uses the faster of the two GPU engines for each query. Taking the minimum
    of two noisy times biases C4 towards the GPU side, most when the two GPU
    engines are close. The analyzer records which engine was selected in each
    round, so readers can check how often the selection switched.
\end{itemize}

\subsection{Construct validity}
\begin{itemize}[leftmargin=*,itemsep=2pt]
  \item \textbf{Timing scope differs at the result boundary.} The Spark timer
    includes collecting results into Python \texttt{Row} objects; the Polars
    timer stops at a host DataFrame and excludes \texttt{to\_dicts()}. Python
    object conversion is therefore partly inside Spark's timer and outside
    Polars'. Result sets are small (at most 20 rows per baseline query and at
    most 175 rows per complex query in the A100 study's saved outputs), so we
    expect the effect to be small, but it is not measured separately.
  \item \textbf{Float64 normalization.} Decimal money and quantity columns are
    cast to Float64 in every engine. This removes decimal-arithmetic costs and
    any GPU decimal support from the measurement.
  \item \textbf{Median-of-medians metrics.} The suite geometric mean weights
    every query equally regardless of its duration; the legacy sum is dominated
    by long queries. Neither is a once-through suite time, throughput, or
    TPC metric.
  \item \textbf{Strict GPU coverage.} Under \texttt{raise\_on\_fail=True},
    \polarsgpu{} does not time a query it cannot fully run on the GPU; a user with
    default settings would see a CPU fallback instead. The Spark plugin can run
    hybrid plans; our plan rule checks operator selection, not the fraction of
    runtime spent on the GPU. For \polarsgpu{}, the streaming executor's
    single-partition fallback still runs on the GPU and is captured as
    warnings, not counted as a CPU fallback.
  \item \textbf{Final-plan evidence} shows which operators were selected, not
    whether each did nonempty work (for example, an anti-join branch whose
    input is empty).
  \item \textbf{Startup excluded.} Session start, JVM launch and table
    registration are outside the timer; they are reported separately.
\end{itemize}

\subsection{External validity}
\begin{itemize}[leftmargin=*,itemsep=2pt]
  \item \textbf{One primary host.} All primary results come from one Lambda
    H100 SXM5 instance with \stat{host.vcpus} vCPUs, local NVMe and \stat{host.host_ram_gib}\,GiB RAM. They do not
    describe distributed Spark clusters, object storage, or equal-cost CPU
    instances. A CPU-only host of the same price could have many more cores.
  \item \textbf{GPU choice limited by availability.} GB300 and H200 instances
    were not available from the provider, so the newest rung is a B200. The
    results say nothing about GPUs that are not on the ladder.
  \item \textbf{Earlier A100 study.} The earlier study
    (\stat{a100_h100.study.label}) ran \stat{a100_h100.study.cpu_threads}
    threads, \stat{a100_h100.study.sessions_per_cell} session per cell and
    Polars \stat{a100_h100.study.polars_version} (against
    \stat{a100_h100.study.polars_version_h100} here), so its comparison with
    the H100 (Figure~\ref{fig:f9}) is descriptive only and its CPU engines are
    not compared; the A100 ladder campaign, which runs this design at 26
    threads, replaces it as the like-for-like A100 result.
  \item \textbf{Version-locked CPU baseline.} Polars 1.44.2 is used because no
    released \texttt{cudf-polars} supports 2.0, so the CPU reference in C1 and
    C4 is an older Polars than the current release. In the side study, 2.0.0
    on the CPU was \stat{bridge_summary.both}.
    Where 2.0.0 is faster, C1 and C4 overstate the GPU's advantage over the
    newest CPU Polars, by roughly that ratio. We cannot measure a GPU engine on
    Polars 2.0; the post-hoc E1/E4 (Section~\ref{sec:exploratory}) are not a
    substitute for a pre-registered contrast.
  \item \textbf{Fixed, untuned configurations.} The Spark settings are those
    of the A100 study, not tuned for the H100 or for each scale; the Polars GPU
    engine uses defaults. Tuning could change any contrast, in either
    direction.
  \item \textbf{Derived, uncertified workload.} Twelve selected or custom
    queries on DuckDB-generated data with fixed parameters are not TPC-H and
    do not represent all analytical workloads. Results may not transfer to
    other query shapes, data distributions, string-heavy or skewed data.
  \item \textbf{Software snapshot.} Results describe specific releases
    (Polars 1.44.2, cudf-polars 26.10.2, Spark 3.5.9, plugin 26.08.1) as of
    October 2026; the Polars GPU engine is labelled open beta and changes
    quickly.
  \item \textbf{GPU coverage gaps} remove two complex queries from every
    \polarsgpu{} contrast, so C1 and C3 complex-suite estimates average over a
    different query set than C2 and C4.
\end{itemize}

\subsection{Validity of the ladder comparisons}
\begin{itemize}[leftmargin=*,itemsep=2pt]
  \item \textbf{Each rung changes the whole host.} Between rungs the GPU,
    device memory, CPU model, RAM, CPU--GPU link and driver
    all change (Table~\ref{tab:hosts}). Cross-host ratios describe hosts, not
    GPUs; no cross-host difference is attributed to the GPU generation alone.
  \item \textbf{26 threads on hosts with more vCPUs.} Every host runs the
    engines with 26 threads, so the A100 and GH200 hosts' larger vCPU counts
    (Table~\ref{tab:hosts}) are not used fully. This keeps engine thread counts equal across the ladder but
    understates what the CPU engines could do on those hosts, and with them
    the C4 reference on that host: on the GH200 especially, a CPU engine using
    all of the host's vCPUs could narrow or reverse a GPU advantage. On the A100 and
    GH200 the idle vCPUs can also absorb operating-system, JVM and telemetry
    threads, which on the H100 and B200, with fewer vCPUs, compete with the engines; this
    slightly favours the larger hosts. Each host's result is a 26-thread
    result.
  \item \textbf{Equal threads, unequal cores.} The hosts can differ in
    simultaneous multithreading and core count, as reported by \texttt{lscpu}
    inside each virtual machine (Table~\ref{tab:hosts}). Where a host has two
    hardware threads per core, the 26 engine threads share half as many
    physical cores as on a host with one, so cross-host ratios for the CPU
    engines, and each host's CPU reference in C1, C2 and C4, also reflect the
    number of physical cores in use, not only the CPU model. The topology is
    guest-reported; how virtual CPUs map to host cores is not visible to us.
  \item \textbf{Arm versus x86 CPU engines.} On the GH200 the CPU engines run
    aarch64 builds of the same versions (Polars wheel, JDK and JIT, Spark, and
    the plugin's aarch64 artifact) on a Grace CPU. Differences in the CPU
    engines across that rung mix the instruction set, compiler output, SIMD
    code paths, memory system and core count of the CPU; they are not a
    measure of the H100 versus the GH200's GPU. Because C4 and C1/C2 use the
    host's own CPU engines as the reference, a change in those contrasts on
    the GH200 can come from the CPU side. The platform-specific artifacts are
    recorded with their SHA-256.
  \item \textbf{Drivers differ.} Each host runs the NVIDIA driver of its
    provider image (Table~\ref{tab:hosts}); drivers are recorded, not
    equalized. The pinned packages are identical, but driver-level behaviour
    (for example, memory management and clock policy) is confounded with the
    host.
  \item \textbf{Each additional host is conditional.} A host's campaign runs
    only if the provider has capacity and the pinned stack passes preflight
    and input verification on that GPU and instruction set; versions are never
    changed to make a host run. The ladder that ran may therefore be
    incomplete (hosts not run: \stat{hosts.not_run_list}). These conditions
    were fixed before any timing and do not depend on results, so a missing
    rung does not select hosts by outcome, but it does limit which
    generations the paper can speak to.
  \item \textbf{Campaigns run at different times.} The hosts' campaigns run
    one after another with identical software, so they share the software
    snapshot but not the moment; provider-side conditions (for example, a
    different physical host behind the same instance type) are not
    controlled.
  \item \textbf{Device and host memory.} The A100's 40\,GB device memory makes
    resource-exhaustion coverage outcomes at SF300 more likely there, which
    shrinks the set of queries in its \polarsgpu{} and \sparkgpu{} contrasts;
    its host RAM, the smallest in the provider's listing
    (Table~\ref{tab:hosts}), also leaves less room for the page cache.
    Cross-host suite ratios are computed over queries
    supported on both hosts, so they can average over different query sets
    for different hosts.
% BEGIN-RAN:b200
\ifbtwohundredran
  \item \textbf{Known Blackwell compiler defect (amendment A4).} On
    Blackwell, the CUDA 12.9 \texttt{ptxas} compiler can miscompile
    \texttt{thrust::exclusive\_scan\_by\_key}, silently corrupting results from
    the second tile onward~\citep{cccl11167}. The reports reproduce it only on
    another Blackwell variant (sm\_120); whether it affects the B200 (sm\_100)
    is not established either way. Reported manifestations include Parquet
    reads of files with many string pages and splits of many buffers, which
    are most likely at SF100 and SF300. The \polarsgpu{} build contains cuDF's
    workaround~\citep{cudf23973}, though other code built by the same compiler
    could still be affected; the \sparkgpu{} plugin JAR predates the
    workaround. The guard is the output-agreement gate: every B200 output is
    compared with \polarscpu{} (Section~\ref{sec:gates}), so a corrupted
    result that changes a query answer fails the round as an integrity failure
    and is never reported as a timing. In the B200 campaign,
    \stat{b200.agreement_cells.polars_gpu.within_gate} of
    \stat{b200.agreement_cells.polars_gpu.total} \polarsgpu{} and
    \stat{b200.agreement_cells.spark_gpu.within_gate} of
    \stat{b200.agreement_cells.spark_gpu.total} \sparkgpu{} suite--scale cells
    had no output outside the gate's tolerance. The largest relative
    differences from \polarscpu{} were
    \stat{b200.agreement_max.polars_gpu.max_rel_diff} (\polarsgpu{}) and
    \stat{b200.agreement_max.spark_gpu.max_rel_diff} (\sparkgpu{}), against
    \stat{agreement_max.polars_gpu.max_rel_diff} and
    \stat{agreement_max.spark_gpu.max_rel_diff} on the H100 (ratios of the
    maxima \stat{xagree.b200.polars_gpu.max_ratio} and
    \stat{xagree.b200.spark_gpu.max_ratio}); the B200's largest difference
    exceeded the H100's in \stat{xagree.b200.polars_gpu.cells_larger} of
    \stat{xagree.b200.polars_gpu.cells} and
    \stat{xagree.b200.spark_gpu.cells_larger} of
    \stat{xagree.b200.spark_gpu.cells} cells, respectively. This is absence of
    evidence that the defect affected these outputs, not proof that it did
    not.
% NOTE: if any within_gate count is below its total, the B200 round
% failed integrity; rewrite this sentence around that outcome.
    Its diagnosis records whether the
    evidence (engine, operator, scale) fits the defect; the defect is not used
    to reclassify a failure. If B200 rounds fail integrity twice, the B200
    campaign is reported as not completed, without version changes. The gate
    compares final answers only: a defect that slows a query or corrupts an
    intermediate result without changing the answer would not be detected.
\fi
% END-RAN:b200
  \item \textbf{The ladder is not independent replication.} Every campaign
    uses the same author, code, software versions, input recipe and provider,
    so the ladder tests whether the results hold on other hardware, not
    whether another team would reproduce them.
\end{itemize}

% ======================================================================
\section{Reproducibility and Artifact}
\label{sec:artifact}

The repository contains the generators, query definitions (SQL and native
Polars), runner, worker, cleanup helper, analyzers and plotting code; the
locked design document; the per-round saved results (every trial time, every
materialized output, Spark final plans, Polars GPU session guard records,
telemetry and manifests); the input manifests; and checksums of all copied
results. Generated Parquet data and NVIDIA binaries are not redistributed; the
generation recipe and pinned dependency versions (with the plugin JAR's
SHA-256) are provided; for the GH200, the aarch64 artifacts and their
SHA-256 are recorded as well. Each host's campaign (A100, H100, GH200, B200,
where run) has its own campaign root, host record and analysis output and is
analyzed separately; the cross-host comparison is a separate output computed
from the per-host results. Table~\ref{tab:jars} identifies the Spark plugin
JAR each host ran: the GH200 uses the plugin's aarch64 (\texttt{arm64}) build
of the same version.

\begin{table}[h]
\centering
\caption{Spark plugin JAR per host, as recorded in each campaign's host
identity (first 16 hexadecimal digits of the SHA-256; the full digest is in
each campaign's \texttt{stats.json}). \hostnotrun{}: campaign not run.}
\label{tab:jars}
\scriptsize
\begin{tabular}{@{}lll@{}}
\toprule
Host & JAR & SHA-256 (prefix) \\
\midrule
A100  & \hoststat{a100}{host.rapids_jar}  & \texttt{\hoststat{a100}{host.rapids_jar_sha256_short}} \\
H100  & \stat{host.rapids_jar}             & \texttt{\stat{host.rapids_jar_sha256_short}} \\
GH200 & \hoststat{gh200}{host.rapids_jar} & \texttt{\hoststat{gh200}{host.rapids_jar_sha256_short}} \\
B200  & \hoststat{b200}{host.rapids_jar}  & \texttt{\hoststat{b200}{host.rapids_jar_sha256_short}} \\
\bottomrule
\end{tabular}
\end{table} The statistics and every number in this paper are
regenerated from saved results by the campaign analyzer, and the paper's
numbers are inserted from its \texttt{stats.json} output rather than typed.
% TODO(author): repository URL, archive DOI (e.g. Zenodo), licence of
% original contributions (unresolved; see NOTICES.md), and TPC/DuckDB notice
% material to accompany the artifact.
Artifact location: \pending{repository URL and archive DOI}.

% ======================================================================
\section{Conclusion}
\label{sec:conclusion}

We designed and pre-registered a four-engine, four-round randomized comparison
of CPU and GPU variants of Polars and Spark on one H100 host, with integrity
gates that do not depend on timings and interval estimates for every contrast.
On the H100, \polarsgpu{} was slower than Polars~1.44.2 on the CPU at
\stat{c1.baseline.summary.below} and faster \stat{c1.both.above_text}; the best
GPU engine
was faster than Polars~1.44.2 on the CPU from
\stat{crossover.baseline.c4.scale} (baseline) and
\stat{crossover.complex.c4.scale} (complex), by
\ratiotext{c4.baseline.sf300} and \ratiotext{c4.complex.sf300} at SF300.
\polarsgpu{} was the faster GPU engine everywhere but rejected
\stat{coverage.polars_gpu.complex.unsupported} of the twelve queries as
unsupported (CPU fallback disabled by design). Polars~2.0.0 on the CPU, which the GPU engine cannot use, was
\stat{bridge_summary.both}, and against it the GPU's advantage was smaller
(exploratory, post-hoc).
% BEGIN-RAN:any
\ifladderran
Repeating the campaign on \stat{hosts.other_ran_list} did not change that
answer qualitatively: the GPU engines were slower on the A100 and not
consistently faster on the GH200 or B200 than on the H100.
\fi
% END-RAN:any
These results describe specific hosts at 26 engine threads, one derived
workload and specific software releases; they do not establish a universal
winner.

\section*{Acknowledgements}
% TODO(author)
\pending{funding, compute provider, reviewers.} No endorsement by NVIDIA,
the Polars project, DuckDB, the Apache Software Foundation or TPC is implied.

\bibliographystyle{plainnat}
\bibliography{references}

\appendix
\section{Per-query speedups at every scale}
\label{app:forest}
Figure~\ref{fig:f1} shows SF300. The remaining scales follow with the same
layout; the complex suite starts at SF30. \chartnote

\begin{figure}[h]
\centering
\begin{subfigure}[t]{0.49\linewidth}\centering\studyfig{f01_forest_sf1}{\linewidth}\caption{SF1}\end{subfigure}\hfill
\begin{subfigure}[t]{0.49\linewidth}\centering\studyfig{f01_forest_sf10}{\linewidth}\caption{SF10}\end{subfigure}\\[4pt]
\begin{subfigure}[t]{0.49\linewidth}\centering\studyfig{f01_forest_sf30}{\linewidth}\caption{SF30}\end{subfigure}\hfill
\begin{subfigure}[t]{0.49\linewidth}\centering\studyfig{f01_forest_sf100}{\linewidth}\caption{SF100}\end{subfigure}
\caption{Per-query speedup ratios with 95\% CIs (four rounds, df~=~3) at SF1
to SF100. \chartnote}
\label{fig:f1-all}
\end{figure}

% BEGIN-RAN:any
\ifladderran
\section{Full contrast tables for the other ladder hosts}
\label{app:ladder}
Each table repeats Table~\ref{tab:contrasts} for one additional host,
computed from that host's campaign alone (four rounds, df~=~3). Hosts whose
campaign did not run are omitted.
% BEGIN-RAN:a100
\ifaonehundredran
\begin{table}[h]
\centering
\caption{A100 campaign: suite-level geometric-mean speedup ratios with
95\% CIs (compare Table~\ref{tab:contrasts}). \chartnote}
\label{tab:contrasts-a100}
\scriptsize
\setlength{\tabcolsep}{4pt}
\begin{tabular}{@{}llcccc@{}}
\toprule
Suite & Scale & \shortstack{C1\\\polarscpu/\\\polarsgpu} & \shortstack{C2\\\sparkcpu/\\\sparkgpu} & \shortstack{C3\\\sparkgpu/\\\polarsgpu} & \shortstack{C4\\\polarscpu/\\best GPU} \\
\midrule
baseline & SF1   & \ratioci{a100.c1.baseline.sf1} & \ratioci{a100.c2.baseline.sf1} & \ratioci{a100.c3.baseline.sf1} & \ratioci{a100.c4.baseline.sf1} \\
         & SF10  & \ratioci{a100.c1.baseline.sf10} & \ratioci{a100.c2.baseline.sf10} & \ratioci{a100.c3.baseline.sf10} & \ratioci{a100.c4.baseline.sf10} \\
         & SF30  & \ratioci{a100.c1.baseline.sf30} & \ratioci{a100.c2.baseline.sf30} & \ratioci{a100.c3.baseline.sf30} & \ratioci{a100.c4.baseline.sf30} \\
         & SF100 & \ratioci{a100.c1.baseline.sf100} & \ratioci{a100.c2.baseline.sf100} & \ratioci{a100.c3.baseline.sf100} & \ratioci{a100.c4.baseline.sf100} \\
         & SF300 & \ratioci{a100.c1.baseline.sf300} & \ratioci{a100.c2.baseline.sf300} & \ratioci{a100.c3.baseline.sf300} & \ratioci{a100.c4.baseline.sf300} \\
\midrule
complex  & SF30  & \ratioci{a100.c1.complex.sf30} & \ratioci{a100.c2.complex.sf30} & \ratioci{a100.c3.complex.sf30} & \ratioci{a100.c4.complex.sf30} \\
         & SF100 & \ratioci{a100.c1.complex.sf100} & \ratioci{a100.c2.complex.sf100} & \ratioci{a100.c3.complex.sf100} & \ratioci{a100.c4.complex.sf100} \\
         & SF300 & \ratioci{a100.c1.complex.sf300} & \ratioci{a100.c2.complex.sf300} & \ratioci{a100.c3.complex.sf300} & \ratioci{a100.c4.complex.sf300} \\
\bottomrule
\end{tabular}
\end{table}
\fi
% END-RAN:a100
% BEGIN-RAN:gh200
\ifghtwohundredran
\begin{table}[h]
\centering
\caption{GH200 campaign: suite-level geometric-mean speedup ratios with
95\% CIs (compare Table~\ref{tab:contrasts}). \chartnote}
\label{tab:contrasts-gh200}
\scriptsize
\setlength{\tabcolsep}{4pt}
\begin{tabular}{@{}llcccc@{}}
\toprule
Suite & Scale & \shortstack{C1\\\polarscpu/\\\polarsgpu} & \shortstack{C2\\\sparkcpu/\\\sparkgpu} & \shortstack{C3\\\sparkgpu/\\\polarsgpu} & \shortstack{C4\\\polarscpu/\\best GPU} \\
\midrule
baseline & SF1   & \ratioci{gh200.c1.baseline.sf1} & \ratioci{gh200.c2.baseline.sf1} & \ratioci{gh200.c3.baseline.sf1} & \ratioci{gh200.c4.baseline.sf1} \\
         & SF10  & \ratioci{gh200.c1.baseline.sf10} & \ratioci{gh200.c2.baseline.sf10} & \ratioci{gh200.c3.baseline.sf10} & \ratioci{gh200.c4.baseline.sf10} \\
         & SF30  & \ratioci{gh200.c1.baseline.sf30} & \ratioci{gh200.c2.baseline.sf30} & \ratioci{gh200.c3.baseline.sf30} & \ratioci{gh200.c4.baseline.sf30} \\
         & SF100 & \ratioci{gh200.c1.baseline.sf100} & \ratioci{gh200.c2.baseline.sf100} & \ratioci{gh200.c3.baseline.sf100} & \ratioci{gh200.c4.baseline.sf100} \\
         & SF300 & \ratioci{gh200.c1.baseline.sf300} & \ratioci{gh200.c2.baseline.sf300} & \ratioci{gh200.c3.baseline.sf300} & \ratioci{gh200.c4.baseline.sf300} \\
\midrule
complex  & SF30  & \ratioci{gh200.c1.complex.sf30} & \ratioci{gh200.c2.complex.sf30} & \ratioci{gh200.c3.complex.sf30} & \ratioci{gh200.c4.complex.sf30} \\
         & SF100 & \ratioci{gh200.c1.complex.sf100} & \ratioci{gh200.c2.complex.sf100} & \ratioci{gh200.c3.complex.sf100} & \ratioci{gh200.c4.complex.sf100} \\
         & SF300 & \ratioci{gh200.c1.complex.sf300} & \ratioci{gh200.c2.complex.sf300} & \ratioci{gh200.c3.complex.sf300} & \ratioci{gh200.c4.complex.sf300} \\
\bottomrule
\end{tabular}
\end{table}
\fi
% END-RAN:gh200
% BEGIN-RAN:b200
\ifbtwohundredran
\begin{table}[h]
\centering
\caption{B200 campaign: suite-level geometric-mean speedup ratios with
95\% CIs (compare Table~\ref{tab:contrasts}). \chartnote}
\label{tab:contrasts-b200}
\scriptsize
\setlength{\tabcolsep}{4pt}
\begin{tabular}{@{}llcccc@{}}
\toprule
Suite & Scale & \shortstack{C1\\\polarscpu/\\\polarsgpu} & \shortstack{C2\\\sparkcpu/\\\sparkgpu} & \shortstack{C3\\\sparkgpu/\\\polarsgpu} & \shortstack{C4\\\polarscpu/\\best GPU} \\
\midrule
baseline & SF1   & \ratioci{b200.c1.baseline.sf1} & \ratioci{b200.c2.baseline.sf1} & \ratioci{b200.c3.baseline.sf1} & \ratioci{b200.c4.baseline.sf1} \\
         & SF10  & \ratioci{b200.c1.baseline.sf10} & \ratioci{b200.c2.baseline.sf10} & \ratioci{b200.c3.baseline.sf10} & \ratioci{b200.c4.baseline.sf10} \\
         & SF30  & \ratioci{b200.c1.baseline.sf30} & \ratioci{b200.c2.baseline.sf30} & \ratioci{b200.c3.baseline.sf30} & \ratioci{b200.c4.baseline.sf30} \\
         & SF100 & \ratioci{b200.c1.baseline.sf100} & \ratioci{b200.c2.baseline.sf100} & \ratioci{b200.c3.baseline.sf100} & \ratioci{b200.c4.baseline.sf100} \\
         & SF300 & \ratioci{b200.c1.baseline.sf300} & \ratioci{b200.c2.baseline.sf300} & \ratioci{b200.c3.baseline.sf300} & \ratioci{b200.c4.baseline.sf300} \\
\midrule
complex  & SF30  & \ratioci{b200.c1.complex.sf30} & \ratioci{b200.c2.complex.sf30} & \ratioci{b200.c3.complex.sf30} & \ratioci{b200.c4.complex.sf30} \\
         & SF100 & \ratioci{b200.c1.complex.sf100} & \ratioci{b200.c2.complex.sf100} & \ratioci{b200.c3.complex.sf100} & \ratioci{b200.c4.complex.sf100} \\
         & SF300 & \ratioci{b200.c1.complex.sf300} & \ratioci{b200.c2.complex.sf300} & \ratioci{b200.c3.complex.sf300} & \ratioci{b200.c4.complex.sf300} \\
\bottomrule
\end{tabular}
\end{table}
\fi
% END-RAN:b200
\fi
% END-RAN:any

\end{document}
